% !TEX program = xelatex
\documentclass[11pt]{ctexart}
%=====================================================================
% theme-violet.tex —— 讲义精讲 · 紫花主题
% 用法：\documentclass[11pt]{ctexart}
%       \input{assets/theme-violet}
%       \begin{document} ... \end{document}
% 只用 ctexart 自带能力 + xcolor/enumitem/fancyhdr/pifont/amsmath，
% 不依赖 tikz/tcolorbox，保证 MiKTeX 可直接编译。
%=====================================================================

\usepackage[a4paper,left=2.0cm,right=2.0cm,top=2.1cm,bottom=2.3cm,%
            headsep=12pt,footskip=22pt]{geometry}
\usepackage{amsmath,amssymb}
\usepackage{xcolor}
\usepackage{array}
\usepackage{enumitem}
\usepackage{pifont}
\usepackage{fancyhdr}
\usepackage[hidelinks]{hyperref}

%---------------------------------------------------------------- 配色
\definecolor{violet}{RGB}{107,78,143}      % 主色：紫
\definecolor{violetdeep}{RGB}{63,44,92}    % 深紫：标题/强调
\definecolor{petal}{RGB}{183,159,208}      % 花瓣：装饰线、花饰
\definecolor{bandviolet}{RGB}{247,243,251} % 浅紫底：档案/提示框
\definecolor{bandcream}{RGB}{252,249,242}  % 米白底：易错/注意框
\definecolor{rulepurple}{RGB}{217,207,230} % 分隔线
\definecolor{inkgray}{RGB}{92,96,110}      % 次要文字

%---------------------------------------------------------------- 字体
% 有 Noto CJK 用 Noto，没有就退回 ctex 默认（Windows 下为宋体/黑体）
\IfFontExistsTF{Noto Serif CJK SC}{\setCJKmainfont{Noto Serif CJK SC}}{}
\IfFontExistsTF{Noto Sans CJK SC}{\setCJKsansfont{Noto Sans CJK SC}}{}

%---------------------------------------------------------------- 版式
\setlist{nosep,leftmargin=1.9em}
\linespread{1.28}
\setlength{\parskip}{3.5pt}
\setlength{\parindent}{0pt}
\setlength{\tabcolsep}{5pt}
\renewcommand{\arraystretch}{1.28}

%---------------------------------------------------------------- 花饰
% \ding{95} = 五瓣小花，\ding{96} = 星形花；都来自 pifont，必定存在
\newcommand{\orn}{\textcolor{petal}{\ding{95}}}          % 行内小标
\newcommand{\ornstar}{\textcolor{petal}{\ding{96}}}      % 星形花
\newcommand{\flowerrule}{%                                % 居中花饰分隔线
  \par\vspace{7pt}%
  \noindent\makebox[\linewidth]{\color{petal}%
    \rule{0.20\linewidth}{0.6pt}\hspace{8pt}{\small\ding{95}}\hspace{4pt}%
    \raisebox{1.1pt}{\tiny\ding{95}}\hspace{4pt}{\small\ding{95}}\hspace{8pt}%
    \rule{0.20\linewidth}{0.6pt}}%
  \par\vspace{7pt}}

%---------------------------------------------------------------- 标题与节头
\newcommand{\doctitle}[3]{%   #1 主标题  #2 副标题  #3 元信息（可用 \\ 分行）
  \begin{center}
    {\color{petal}\small \orn\enspace #2\enspace \ornstar}\\[7pt]
    {\LARGE\bfseries\color{violetdeep} #1}\\[10pt]
    {\color{petal}\rule{0.55\linewidth}{1pt}}\\[4pt]
    {\small\color{inkgray} #3}
  \end{center}\vspace{2pt}}

\newcommand{\hdr}[2]{%   #1 章节号（如 §0）  #2 标题
  \par\vspace{12pt}%
  \noindent\colorbox{violetdeep}{\color{white}\bfseries\small\;#1\;}\hspace{7pt}%
  {\color{violetdeep}\large\bfseries #2}\par\nobreak
  \vspace{3pt}{\color{petal}\hrule height 1.2pt}\vspace{6pt}}

\newcommand{\subhdr}[1]{%
  \par\vspace{6pt}%
  \noindent{\color{violet}\bfseries #1}\par\vspace{1pt}}

%---------------------------------------------------------------- 提示框
\newcommand{\callout}[1]{% 浅紫底：讲义档案、读法说明
  \par\medskip\noindent
  \fcolorbox{rulepurple}{bandviolet}{\parbox{\dimexpr\linewidth-2\fboxsep-2\fboxrule}{\vspace{2pt}#1\vspace{2pt}}}%
  \par\medskip}

\newcommand{\warnbox}[1]{% 米白底：易错、边界条件
  \par\medskip\noindent
  \fcolorbox{rulepurple}{bandcream}{\parbox{\dimexpr\linewidth-2\fboxsep-2\fboxrule}{\vspace{2pt}#1\vspace{2pt}}}%
  \par\medskip}

\newcommand{\gapitem}[2]{% #1 缺口类型（紫） #2 内容
  \par\vspace{2pt}\noindent
  \hangindent=1.4em\hangafter=1
  {\color{petal}\small\ding{96}}\enspace{\bfseries\color{violetdeep}#1}\enspace #2\par}

% 前提旁注：只放在「需要理解的那一步」后面，给陈述 + 本题代入，不重证。
\newcommand{\aside}[1]{%
  \par\vspace{1pt}\noindent
  \hangindent=1.6em\hangafter=1
  {\small\color{inkgray}%
    \hspace*{0.6em}{\color{petal}$\triangleright$}\hspace{2pt}#1}%
  \par\vspace{2pt}}

%---------------------------------------------------------------- 表格
% 本机 MiKTeX 的 colortbl 与 array 版本不匹配（\rowcolor 会报错），
% 故表格统一用「上粗紫线 + 表头 + 中细线」的三线式，不依赖 colortbl。
\newcommand{\tabtop}{\par\vspace{3pt}{\color{violet}\hrule height 1.2pt}\vspace{1pt}}
\newcommand{\tabmid}{\par\vspace{1pt}{\color{petal}\hrule height 0.5pt}\vspace{1pt}}
\newcommand{\tabbot}{\par\vspace{3pt}{\color{rulepurple}\hrule height 0.5pt}\vspace{8pt}}
\newcommand{\thead}[1]{{\bfseries\color{violetdeep}#1}}

%---------------------------------------------------------------- 页眉页脚
\pagestyle{fancy}
\fancyhf{}
\renewcommand{\headrulewidth}{0.4pt}
\renewcommand{\headrule}{\hbox to\headwidth{\color{rulepurple}\leaders\hrule height \headrulewidth\hfill}}
\fancyhead[L]{\small\color{inkgray}\songti \doctag}
\fancyhead[R]{\small\color{inkgray}\songti \coursetag}
\fancyfoot[C]{\small\color{inkgray}\orn\enspace{\color{violet}\thepage}\enspace\orn}

% 每篇讲义在导言区设这两个值
\newcommand{\doctag}{}
\newcommand{\coursetag}{}
\newcommand{\settags}[2]{\renewcommand{\doctag}{#1}\renewcommand{\coursetag}{#2}}

%---------------------------------------------------------------- 数学记号
\newcommand{\vect}[1]{\boldsymbol{#1}}
\newcommand{\dd}{\,\mathrm{d}}
\newcommand{\ddx}{\frac{\mathrm{d}}{\mathrm{d}x}}
\addtolength{\jot}{2pt}



\settags{C3-3 · 精讲}{电工电子技术}
% —— 本节只在本文件内改写两个节头宏，使标题不与正文分家（不修改主题文件）——
\renewcommand{\hdr}[2]{%
  \par\vspace{12pt}%
  \noindent\colorbox{violetdeep}{\color{white}\bfseries\small\;#1\;}\hspace{7pt}%
  {\color{violetdeep}\large\bfseries #2}\par\nobreak
  \vspace{3pt}{\color{petal}\hrule height 1.2pt}\nopagebreak[4]\vspace{6pt}}
\renewcommand{\subhdr}[1]{%
  \par\vspace{6pt}%
  \noindent{\color{violet}\bfseries #1}\par\nopagebreak[4]\vspace{1pt}}


\begin{document}

\doctitle{3.5--3.6\quad 滤波与稳压：从纹波到稳压值}
{稳压管工作点怎么判、限流电阻怎么取、三端稳压器怎么认管脚}
{依据：课件 Chapter Three-3（p1--p38，共 38 页；本讲义页码一律指本册 PDF 物理页，引用他册时写成 [C3-2 p44] 的形式）\ \textperiodcentered\ 2026-10-23}

%=============================================================
\hdr{§0}{课件坐标、记号与三处拼接}

\subhdr{0.1 这 38 页讲了哪几块，各在哪几页}

课件不是一条直线，而是「复习两块、正文三块」：开头七页把上一册的整流结论与两道选择题重放一遍，
p8 起进入 3.5 滤波，p15 起进入 3.6 稳压，最后两页才讲集成稳压器。中间 p28--p36 是同一个例题
（Example 2）的九帧连续画面，内容几乎不增，只多出老师的手写笔画——读的时候当成一页即可。

\par\vspace{2pt}
\noindent\begin{tabular}{@{}p{4.0cm}p{2.0cm}p{10.0cm}@{}}
\tabtop
\thead{课件小节} & \thead{页码} & \thead{本讲义对应（补了什么）}
\\
\tabmid
Review 与两道选择题 & p1--p5 & §1.3（答案与理由，课件只有题没有解）
\\
Design / 灯珠实物 & p6--p7 & §1.4（把整条链算一遍，含 $16.8\,\mathrm{V}$ 的核算）
\\
3.5.2 电容滤波 E1 & p8--p11 & §2.1--§2.4（四问全程，并点出 p9 的 $34.5$ 笔误）
\\
3.5.3 电感滤波 & p12 & §2.5（为什么「低压大电流」用它，纹波分压的定量表）
\\
3.5.4 复式滤波 & p13 & §2.6（$\pi$ 型的命名、RC 版的代价与画图不一致）
\\
3.6 路标（实物） & p14 & §2.7（两件实物各对应后文哪一节）
\\
3.6.1 稳压管器件 & p15--p18 & §3（击穿为什么不等于损坏、四个参数怎么用）
\\
3.6.1 稳压电路与过程 & p19--p20 & §4（把箭头链补成三个方程与两条定量式）
\\
Example 1 (a)--(g) & p21--p27 & §6（七问全程）；判据与限流电阻取法在 §5
\\
Example 2（九帧） & p28--p36 & §8.1--§8.3（读数 $000$ 的分支诊断，课件没给答案）
\\
3.6.2 集成稳压器 & p37--p38 & §8.4--§8.6（管脚、XX 的含义、压差「[教材补充]」）
\tabbot
\end{tabular}

\aside{与技能默认骨架（§2--§3 两条主线、§4 讲等价）不同，本讲义按\textbf{课件自己的链}排：
滤波在 §2、稳压在 §3--§6，因为这份课件的顺序就是「先把脉动削平、再把数值钉住」，
把稳压提到前面反而要先解释一堆「$\sqrt2U_2$ 是哪来的」。§0--§12 的编号与技能一致，只是节内内容按链重排。}

\subhdr{0.2 记号：先统一，后文不再回头}

\par\vspace{1pt}{\small
\noindent\begin{tabular}{@{}p{2.5cm}p{11.3cm}@{}}
\tabtop
\thead{记号} & \thead{约定（本册一律按此，不与他册混）}
\\
\tabmid
$U_2$ & 变压器\textbf{二次电压有效值}（交流量；峰值写 $U_{\mathrm{m}}=\sqrt2U_2$）
\\
$U_{\mathrm{I}}$ & 稳压环节的\textbf{输入}电压平均值，也就是滤波电容两端那个「还没稳」的直流（p19 写作 $U_{\mathrm{I}}$）
\\
$U_{\mathrm{O}},\ I_{\mathrm{O}}$ & 负载上的\textbf{输出}电压、电流平均值（p9 写作 $U_{\mathrm{O}}$、$I_{\mathrm{O}}$）
\\
$U_{\mathrm{Z}}$ & 稳压管的\textbf{稳压值}，取正值（反向击穿电压的绝对值，p18）
\\
$I_{\mathrm{Z}}$ & 流过稳压管的\textbf{反向}电流，取正值（p20 图中的 $I_{\mathrm{Z}}$ 箭头）
\\
$I_{\mathrm{Zmin}},\ I_{\mathrm{Zmax}}$ & 稳压管的\textbf{最小}、\textbf{最大}工作电流（p18 用两段文字定性给出）
\\
$r_{\mathrm{Z}}$ & 稳压管击穿段的\textbf{动态电阻} $r_{\mathrm{Z}}=\Delta U/\Delta I$（p18）
\\
$R,\ R_{\mathrm{L}}$ & 限流（串联）电阻、负载电阻；本册的 $R$ 一律指稳压管前面那只限流电阻
\\
$U_{\mathrm{F}}$ & 二极管\textbf{正向压降}，Example 1 取 $0.5\,\mathrm{V}$（p21）
\tabbot
\end{tabular}}

\warnbox{\textbf{最容易混的两个 $U$}：$U_2$ 是\textbf{交流有效值}，$U_{\mathrm{O}}$、$U_{\mathrm{I}}$ 是\textbf{直流平均值}。
$U_{\mathrm{O}}=1.2U_2$ 这条式子的左右两边量纲相同但物理意义不同，
所以要写「$U_2=U_{\mathrm{O}}/1.2$ 反求」时，必须清楚自己是在用一条\textbf{经验系数}换算，不是在同一本账上加减。}

\subhdr{0.3 三处拼接（课件自己的缝）}

\gapitem{拼接缺口}{p11 的两个极限（负载开路 $35.4\,\mathrm{V}$、电容支路断开 $22.5\,\mathrm{V}$）
已经把「滤波不能稳」说透了，但课件\textbf{没有一句话}把这件事接到 p12 之后的 3.6 上。
中间缺的那句是：滤波决定了 $U_{\mathrm{I}}$ 的\textbf{大小和纹波}，但 $U_{\mathrm{I}}$ 仍随电网与负载漂；
稳压级要做的不是「再削一次纹波」，而是\textbf{把 $U_{\mathrm{O}}$ 钉在一个与 $U_{\mathrm{I}}$ 无关的数上}。
本讲义 §1.1 补这句，§2.4 与 §4 各再落一次。}

\gapitem{拼接缺口}{p15--p18 讲器件的四个参数，p19 突然给整机电路，p21 直接进例题。
「参数怎么变成判据」这一步（$I_{\mathrm{Z}}$ 必须落在 $I_{\mathrm{Zmin}}$ 与 $I_{\mathrm{Zmax}}$ 之间如何\textbf{验算}）
课件从未写成可执行的步骤。本讲义 §5 给出三步法，§7.2 用 Example 1(g) 演示它翻车的现场。}

\gapitem{版面重复}{p28--p36 是 Example 2 的九帧，同一张底图只换手写笔画；p23 与 p24 两个子图
(c)、(d) 的端子极性标注在位图上完全一致（输入上「$-$」下「$+$」、输出上「$+$」下「$-$」），
两者差异在扫描图与文字层里都读不出来。§0 的坐标表按去重后的内容计，§10 记两条，涉及 (c)(d) 处标 [待确认]。}

\callout{\textbf{本讲义怎么用}：\\
第一遍只看 §1.1 的链条、§5 的三步法、§6 的答案表、§9 的公式卡——这四处直接决定题做不做得对；\\
第二遍看 §3、§4、§7——这三处解释「为什么必须这么判」，是换一个电路还能不能算对的关键；\\
\textbf{§10 缺口清单}逐条对应课件某页某句：课件哪句话让你卡住，就去那里找补在哪。}

%=============================================================
\hdr{§1}{两条链：本册到底在解决什么}

\subhdr{1.1 整流--滤波--稳压：每一段只解决一件事}

直流稳压电源是四段串联，每一段只改一个属性，记住这一句就不会把公式用串：

\par\vspace{2pt}
\noindent\begin{tabular}{@{}p{2.6cm}p{4.3cm}p{8.9cm}@{}}
\tabtop
\thead{环节} & \thead{改的是} & \thead{留下什么问题}
\\
\tabmid
变压器 & 把 $220\,\mathrm{V}$ 降到 $U_2$（$\sim$ 号那两圈） & 还是交流，仍有负半周
\\
整流 & \textbf{方向}：把交流变成单向脉动 & 方向对了，但数值在 $0$ 与 $\sqrt2U_2$ 之间大幅波动
\\
滤波 & \textbf{脉动}：用电容或电感把波动削平 & 平均值抬高了，但\textbf{抬到多少取决于负载}（p11）
\\
稳压 & \textbf{数值}：把 $U_{\mathrm{O}}$ 钉在一个固定值上 & 稳压管要吃掉多余电压，于是有了电流限额问题
\tabbot
\end{tabular}

一句话：\textbf{整流管方向、滤波管纹波、稳压管数值}。三件事互不替代，
因此 p11 的 $35.4\,\mathrm{V}$／$22.5\,\mathrm{V}$ 与 p19 的稳压级不是重复劳动，而是两笔不同的账。

\subhdr{1.2 p11 的两个极限就是「为什么还需要稳压」的证据}

p11 在同一个电路上问了两件事（第（3）问负载开路、第（4）问电容支路断开），数值分别是
\[
U_{\mathrm{O}}\big|_{\text{开路}}=\sqrt2U_2=35.4\,\mathrm{V},
\qquad
U_{\mathrm{O}}\big|_{\text{断容}}=0.9U_2=22.5\,\mathrm{V}.
\]
同一个电源，负载一轻电压就往上跑到 $35.4\,\mathrm{V}$，电容一坏就掉到 $22.5\,\mathrm{V}$——
\textbf{变化幅度 $57\%$}。这就是 3.6 存在的理由：只靠滤波，$U_{\mathrm{O}}$ 从来不是一个「数」，
而是一个随负载与元件状态漂的区间。稳压级要把这个区间压成一点。

\aside{课件 p11 只给两行答案，没有比较这两行；本讲义的 $57\%$ 由 $(35.4-22.5)/22.5$ 算出。
注意这两个数\textbf{不是同一工况}（一个是负载开路、一个是电容断开），
所以 $57\%$ 只用来提示「变化之大」，不能当作指标写进设计计算。}

\subhdr{1.3 p1 的两道选择题：答案与理由}

课件 p1 只有题干与选项，没有答案。两题的图都在 p1 上半页，答案取决于「查总表哪一列」。

\textbf{第 1 题}：\emph{Find the maximum reverse voltage of the Diode}（选项 $0$、$U_2$、$\sqrt2U_2$、$2\sqrt2U_2$）。
图是\textbf{半波整流 + 电容}（次级经一只二极管接 $R$，$R$ 两端并了电容 $C$）。查 p1 右下角那张总表
「半波 / 有电容」一列：$U_{\mathrm{Rm}}=2\sqrt2U_2$，所以选 \textbf{d}。

为什么有电容会翻倍：正半周电容被充到 $+\sqrt2U_2$，之后阴极被电容钳在这个正峰上；
到负半周，阳极被次级拉到 $-\sqrt2U_2$，管子两端电压是两端峰值之差
\[
u_{\mathrm{D}}=u_{\mathrm{a}}-u_{\mathrm{K}}=(-\sqrt2U_2)-(+\sqrt2U_2)=-2\sqrt2U_2 .
\]
\textbf{一端被电容顶在正峰、另一端被次级拉到负峰，两个峰叠加}——这就是那个额外的 $2$。

\textbf{第 2 题}：\emph{Given that $I_{\mathrm{O}}=50\,\mathrm{mA}$, find $I_{\mathrm{D}}$}（选项 $50$、$25$、$12.5$、$60\,\mathrm{mA}$）。
图是\textbf{桥式}全波。查总表「全波 / 无电容」或「全波 / 有电容」两列：$I_{\mathrm{D}}=\tfrac12I_{\mathrm{O}}$，
所以 $I_{\mathrm{D}}=25\,\mathrm{mA}$，选 \textbf{b}。机制是两只管子分管两个半周，负载电流被分成两半。

两题合起来给了第三个数：p1 下半个题干写 \emph{Given that $U_2=100\,\mathrm{V}$, find the maximum reverse voltage}，
对桥式电路是 $U_{\mathrm{Rm}}=\sqrt2\times100\,\mathrm{V}=141.4\,\mathrm{V}$。
\textbf{同一个 $U_2=100\,\mathrm{V}$，半波带电容要 $282.8\,\mathrm{V}$、桥式只要 $141.4\,\mathrm{V}$}——
这也是桥式被普遍采用的原因之一。

\subhdr{1.4 p7 的两串灯珠：把整条链算一遍}

p7 用两样实物把设计任务落地：上面一套是 $220\,\mathrm{V}$ 交流经变压器降到 $14\,\mathrm{V}$、
桥式整流、$1000\,\mu\mathrm{F}$ 电容滤波，再经 $R_1=20\,\Omega/3\,\mathrm{W}$ 限流点亮三只串联 LED（$350\,\mathrm{mA}$ 标称）；
下面一套是交流经 $R_1\parallel C_1$、桥式整流、$C_2/C_3$ 与 $RV$ 之后驱动一整串 LED。
两套都在说明同一件事：\textbf{限流元件串在电源与灯串之间}，它的阻值决定电流。

把上面那套的数算一遍。滤波后的直流按图上标注的 \emph{$V_{\mathrm{DC}}=1.2V_{\mathrm{in}}$}：
\[
V_{\mathrm{DC}}=1.2\times14\,\mathrm{V}=16.8\,\mathrm{V}.
\]
若三只 LED 每只压降 $3.5\,\mathrm{V}$，灯串共 $10.5\,\mathrm{V}$，落在 $R_1$ 上的电压是
\[
U_{R1}=16.8-10.5=6.3\,\mathrm{V},
\qquad
I=\frac{6.3}{20}=0.315\,\mathrm{A}=315\,\mathrm{mA},
\qquad
P_{R1}=\frac{6.3^2}{20}=1.98\,\mathrm{W}.
\]
$315\,\mathrm{mA}$ 与标称 $350\,\mathrm{mA}$ 相差 $10\%$，$1.98\,\mathrm{W}$ 落在 $3\,\mathrm{W}$ 定额之内——
\textbf{这两个吻合反过来证明：图上的 $3.5\,\mathrm{V}$ 是「一只 LED」的压降，不是整串的}。

\aside{图上只画了一个 $3.5\,\mathrm{V}$ 的竖向箭头，无法从图中判定它跨的是一只还是三只。
本讲义的判定依据是 $R_1$ 的功率定额：若按整串 $3.5\,\mathrm{V}$ 算，$I=(16.8-3.5)/20=665\,\mathrm{mA}$，
$P_{R1}=8.84\,\mathrm{W}$，与印着的「$3\,\mathrm{W}$」差三倍，物理上不成立。这一步是\textbf{核算}，不是课件原文。}

这三行值得单独记，因为它把一个常被忽略的问题摆出来了：\textbf{限流电阻的阻值不是一个设计自由度，
而是被「电源电压 $-$ 负载压降」与「所需电流」两者夹住的}。§5 讲稳压管的限流电阻时，用的是同一种夹法。

%=============================================================
\hdr{§2}{滤波电路（3.5，p8--p14）}

\subhdr{2.1 E1 四问：题面与答案总表}

\callout{\textbf{[E1]}（p8 原文）\quad Bridge rectifier, C filter circuit with a line frequency of 50 Hz,
$R_{\mathrm{L}}=100\,\Omega$, output DC $U_{\mathrm{O}}=30\,\mathrm{V}$. Try to：\\
（1）Select the diodes；（2）Select the capacitor；\\
（3）Find the output voltage when the load is open circuit；\\
（4）Find the output voltage when the capacitor branch is open.}

p8 的右下方配着那张四列总表（同一张表在 p1、p28--p36 各出现一次），它是全册的「查表依据」，
所以先把表抄正一次，后文一律引用它：

\par\vspace{2pt}
\noindent\begin{tabular}{@{}p{3.4cm}p{2.9cm}p{2.9cm}p{2.9cm}p{2.9cm}@{}}
\tabtop
\thead{（3）Summary} & \thead{半波\newline 无电容} & \thead{半波\newline 有电容} & \thead{全波\newline 无电容} & \thead{全波\newline 有电容}
\\
\tabmid
$U_{\mathrm{O}}$ & $0.45U_2$ & $U_2$ & $0.9U_2$ & $1.2U_2$
\\
$I_{\mathrm{O}}$ & $U_{\mathrm{O}}/R_{\mathrm{L}}$ & $U_{\mathrm{O}}/R_{\mathrm{L}}$ & $U_{\mathrm{O}}/R_{\mathrm{L}}$ & $U_{\mathrm{O}}/R_{\mathrm{L}}$
\\
$I_{\mathrm{D}}$ & $I_{\mathrm{O}}$ & $I_{\mathrm{O}}$ & $\tfrac12I_{\mathrm{O}}$ & $\tfrac12I_{\mathrm{O}}$
\\
$U_{\mathrm{Rm}}$ & $\sqrt2U_2$ & $2\sqrt2U_2$ & $\sqrt2U_2$ & $\sqrt2U_2$
\tabbot
\end{tabular}

\aside{「全波 $+$ 有电容」一列的 $U_{\mathrm{Rm}}$ 仍写 $\sqrt2U_2$，与「全波无电容」相同；
半波那一列却从 $\sqrt2U_2$ 翻到 $2\sqrt2U_2$。差别在于\textbf{二次侧有没有一条直流通路}：
半波时负载把二次侧的一端直接钉在负母线上，另一端就被变压器按到 $-\sqrt2U_2$，
而阴极又被电容钳在 $+\sqrt2U_2$，两端相加得 $2\sqrt2U_2$；
桥式的二次侧两端都只经二极管接母线、没有直流通路，两端电位不会各自跑到 $\pm\sqrt2U_2$ 的两个极端，
任一只关断的管子两端最多只差一个 $\sqrt2U_2$。课件只给表值，这段机制是本讲义补的（[教材补充]）。}

解题的第一步不是算，是\textbf{定位}：题面三件事分别锁定一列、一条电流关系、一个反解方向。
「Bridge rectifier」锁定「全波」两列；「C filter」锁定「有电容」那一列；
「$U_{\mathrm{O}}=30\,\mathrm{V}$」是\textbf{平均值}，所以从 $U_{\mathrm{O}}=1.2U_2$ 反求 $U_2$。

\par\vspace{2pt}
\noindent\begin{tabular}{@{}p{2.0cm}p{5.5cm}p{3.3cm}p{3.6cm}@{}}
\tabtop
\thead{量} & \thead{算式} & \thead{本题值} & \thead{出处}
\\
\tabmid
$U_2$ & $U_{\mathrm{O}}/1.2=30/1.2$ & $25\,\mathrm{V}$ & p9
\\
$I_{\mathrm{O}}$ & $U_{\mathrm{O}}/R_{\mathrm{L}}=30/100$ & $0.3\,\mathrm{A}$ & p9
\\
$I_{\mathrm{D}}$ & $I_{\mathrm{O}}/2$ & $0.15\,\mathrm{A}$ & p9
\\
$U_{\mathrm{Rm}}$ & $\sqrt2U_2=\sqrt2\times25$ & $35.4\,\mathrm{V}$ & p9
\\
$I_{\mathrm{F}}$ & $\ge2I_{\mathrm{D}}$ & $\ge0.3\,\mathrm{A}$ & p9
\\
$U_{\mathrm{R}}$ & $\ge U_{\mathrm{Rm}}$ & $\ge35.4\,\mathrm{V}$ & p9（课件误写 $34.5$）
\\
$C$ & $\ge(1.5\sim2.5)/(R_{\mathrm{L}}f)$ & $300\sim500\,\mu\mathrm{F}$ & p10
\\
$U_{\mathrm{CN}}$ & $\ge\sqrt2U_2$ & $\ge35.4\,\mathrm{V}$ & p10
\\
（3）开路 & $\sqrt2U_2$ & $35.4\,\mathrm{V}$ & p11
\\
（4）断容 & $0.9U_2$ & $22.5\,\mathrm{V}$ & p11
\tabbot
\end{tabular}

\subhdr{2.2 （1）选二极管（p9）：四个数，一个笔误}

p9 把四步压缩成四行式子，顺序是 U2、IO、ID、URm、IF、UR：

\begin{enumerate}
\item 反求二次电压 $U_2=U_{\mathrm{O}}/1.2=25\,\mathrm{V}$（这一步决定后面全部数字，见 §2.1 的定位）；
\item 负载电流 $I_{\mathrm{O}}=U_{\mathrm{O}}/R_{\mathrm{L}}=0.3\,\mathrm{A}$；
\item 每管电流 $I_{\mathrm{D}}=\tfrac12I_{\mathrm{O}}=0.15\,\mathrm{A}$（桥式两只管子分管两个半周）；
\item 两个选管条件 $I_{\mathrm{F}}\ge2I_{\mathrm{D}}=0.3\,\mathrm{A}$、$U_{\mathrm{R}}\ge U_{\mathrm{Rm}}=35.4\,\mathrm{V}$。
\end{enumerate}

\warnbox{\textbf{课件笔误（必须点名）}：p9 倒数第二行写 \emph{``$U_{\mathrm{R}}\ge U_{\mathrm{Rm}}=34.5\,\mathrm{V}$''}，
而同一页上一行刚算出 $U_{\mathrm{Rm}}=\sqrt2\times25\,\mathrm{V}=35.4\,\mathrm{V}$。
\textbf{$34.5$ 是 $35.4$ 的数字转置}，两个数在同一页自相矛盾。本讲义一律按 $35.4\,\mathrm{V}$ 走：
$\sqrt2\times25=35.355\ldots\approx35.4$，复算见 §12 的核对说明。
后果不轻——耐压取 $34.5\,\mathrm{V}$ 会把阈值定低，选到 $35\,\mathrm{V}$ 档的管子正好卡在边缘。}

\textbf{为什么 $I_{\mathrm{F}}$ 要乘 $2$，而 $U_{\mathrm{R}}$ 不乘}：$I_{\mathrm{F}}$ 是\textbf{发热}定额，
看的是电流的\textbf{有效值}；电容滤波下二极管只在电容电压低于 $\sqrt2U_2$ 的短促尖峰里导通，
电流是窄脉冲，其有效值约为平均值的 $2$ 倍。$U_{\mathrm{R}}$ 是\textbf{耐压}定额，看的是\textbf{峰值}，
峰值已经由 $\sqrt2U_2$ 算进去了，不再乘。
\aside{课件的 $\sqrt2U_2$ 与 $I_{\mathrm{F}}\ge2I_{\mathrm{D}}$ 都是从别处（p1 的总表、上一册 §6.3）搬来的结论，
p9 本身没有解释这两个「$2$」为什么一个有一个没有。本段的波形因数是补的，属于 [教材补充]。}

课件结论是「See Appendix, Select four diodes of 2CZ53B」／「查附录，选用 2CZ53B 4 个」。
本册没有附录，2CZ53B 的具体额定值无法在本册核对，标 [待确认]；
可核对的是它必须同时满足 $I_{\mathrm{F}}\ge0.3\,\mathrm{A}$ 与 $U_{\mathrm{R}}\ge35.4\,\mathrm{V}$。

\subhdr{2.3 （2）选电容（p10）：经验式的等价写法}

p10 只给了一行式子，来历没有交代：
\[
C\ge\frac{1.5\sim2.5}{R_{\mathrm{L}}f}
=\frac{1.5\sim2.5}{100\times50}\,\mathrm{F}
=(300\sim500)\,\mu\mathrm{F}.
\]
把它换一种写法，物理意义就出来了。令 $T=1/f=20\,\mathrm{ms}$ 为周期，则
\[
C\ge\frac{1.5\sim2.5}{R_{\mathrm{L}}f}
\iff
\tau=R_{\mathrm{L}}C\ge(1.5\sim2.5)T=(30\sim50)\,\mathrm{ms}.
\]
即\textbf{电容的放电时间常数要不小于 $1.5\sim2.5$ 个周期}：电容在每个峰值充一次电，
然后靠这一个 $\tau$ 撑到下一个峰值，撑得越久，掉得越少、纹波越小。
这也是「电容越大越接近峰值」的定量版本——与 §1.2 的 $35.4\,\mathrm{V}$ 是同一个机制的两个方向。

\aside{同一页还给了耐压条件 $U_{\mathrm{CN}}\ge\sqrt2U_2=35.4\,\mathrm{V}$，
所以标称必须取 $50\,\mathrm{V}$ 档而不是 $35\,\mathrm{V}$ 档——$35.4\,\mathrm{V}$ 是\textbf{下限}，没有裕量。}

\subhdr{2.4 （3）（4）两个极限（p11）：同一只电容的两个极端}

\textbf{（3）负载开路}：$R_{\mathrm{L}}\to\infty$，电容没有放电回路，充到峰值就一直保持，
\[
U_{\mathrm{O}}=U_{\mathrm{C}}=\sqrt2U_2=35.4\,\mathrm{V}.
\]
\textbf{（4）电容支路断开}：回到纯桥式整流，查「全波无电容」一列，
\[
U_{\mathrm{O}}=0.9U_2=0.9\times25\,\mathrm{V}=22.5\,\mathrm{V}.
\]
两问合起来说明一件事：\textbf{滤波电容既是「抬高器」也是「不稳定源」}。
它把 $0.9U_2=22.5\,\mathrm{V}$ 抬到 $1.2U_2=30\,\mathrm{V}$（抬高 $33\%$），
同时把「输出随负载变化」这条路打开了（$30\to35.4\,\mathrm{V}$）。
真正的稳压要靠下一级，这正是 p19 之后的内容。

\subhdr{2.5 电感滤波（3.5.3，p12）：为什么「低压大电流」用它}

p12 的电路是桥式整流后在正母线上\textbf{串联}一只电感 $L$（扼流圈），再接负载。
课件只有一句英文结论：\emph{Normally used in the situation when the output voltage is relatively lower
while the current through the load is relatively larger.} 这句话背后的机制是：

电感对直流是短路、对交流是高阻抗，所以它\textbf{不抬高}直流平均值（全波约 $U_{\mathrm{O}}\approx0.9U_2$），
只把交流纹波按分压比压下去。设全波纹波的主频为 $2f=100\,\mathrm{Hz}$，则负载上的纹波是
\[
\frac{\Delta U_{\mathrm{O}}}{\Delta U_{\mathrm{in}}}
=\frac{R_{\mathrm{L}}}{\sqrt{R_{\mathrm{L}}^2+(2\omega L)^2}},
\qquad \omega=2\pi f .
\]
代入 $R_{\mathrm{L}}=100\,\Omega$、$f=50\,\mathrm{Hz}$ 算一张小表：

\par\vspace{2pt}
\noindent\begin{tabular}{@{}p{3.0cm}p{3.2cm}p{3.6cm}p{5.1cm}@{}}
\tabtop
\thead{$L$} & \thead{$2\omega L$} & \thead{纹波传到负载的比例} & \thead{结论}
\\
\tabmid
$0.05\,\mathrm{H}$ & $63\,\Omega$ & $0.954$ & 几乎没滤（$L$ 太小）
\\
$0.1\,\mathrm{H}$ & $126\,\Omega$ & $0.847$ & 只压掉一成半
\\
$0.5\,\mathrm{H}$ & $628\,\Omega$ & $0.303$ & 压到三成
\\
$1\,\mathrm{H}$ & $1257\,\Omega$ & $0.157$ & 压到一成半
\tabbot
\end{tabular}

这张表说明两点：（1）电感滤波要\textbf{有效}，必须 $2\omega L\gg R_{\mathrm{L}}$，
所以它天生适合\textbf{负载电阻小}（即电流大）的场合——$R_{\mathrm{L}}$ 越小，分压比越好看；
（2）它\textbf{不抬高}输出电压，所以适合「$U_2$ 本来就不高、又要大电流」的场合。
这两条合起来正是 p12 那句英文的意思。反过来，小电流场合用电感既贵又没效果，那里用 $\pi$ 型 RC 更划算（§2.6）。

\subhdr{2.6 复式滤波（3.5.4，p13）：$\pi$ 型的命名与 RC 版的代价}

p13 并排给了两种复式滤波：（1）$\pi$ 型 LC、（2）$\pi$ 型 RC。名字来自形状：
\textbf{并联电容 $C_1$ — 串联元件 — 并联电容 $C_2$}，三段画在一起像一个希腊字母 $\pi$。
两种的区别只在中间那只串联元件是 $L$ 还是 $R$。

\textbf{（1）$\pi$ 型 LC}：$C_1$ 先把脉动削一遍，$L$ 再把剩下的交流压住，$C_2$ 再守一道。
串联的电感\textbf{不消耗直流功率}，所以适合电流较大的场合；代价是体积、重量与成本。

\textbf{（2）$\pi$ 型 RC}：把 $L$ 换成 $R$ 便宜得多，但 $R$ 上\textbf{有直流压降}：
全波整流时
\[
U_{\mathrm{O}}\approx0.9U_2-I_{\mathrm{O}}R .
\]
例如 $U_2=20\,\mathrm{V}$、$R=50\,\Omega$、$I_{\mathrm{O}}=20\,\mathrm{mA}$，压降是 $1\,\mathrm{V}$，
输出约 $17\,\mathrm{V}$（无电容时 $18\,\mathrm{V}$），还可以接受；
但若 $I_{\mathrm{O}}=200\,\mathrm{mA}$，同样的 $R$ 会吃掉 $10\,\mathrm{V}$ 并使 $R$ 发热 $2\,\mathrm{W}$——
所以 \textbf{RC 版只适合小电流}，这正是它与 LC 版的分工。

\warnbox{\textbf{课件 p13 的版面不一致（点名）}：(2) 的标题写 \emph{$\pi$-type RC Filter}，
但图里的串联元件仍画成\textbf{线圈符号}并标着 $L$，只在正上方多了一个红字 $R$。
按标题与电路理论，串联元件应是\textbf{电阻 $R$}；本讲义一律按「$R$」讲，
课件该图按 $R$ 理解，画线圈处视为版面残留，标 [待确认]。}

\subhdr{2.7 p14：3.6 的路标}

p14 是一页过渡：左边一张成品电源板实物，右边并列两件器件——
一只玻璃封装的\textbf{稳压二极管}、一只 TO-220 封装的\textbf{三端稳压器}。
课件的正文只有两行 \emph{Content: （1）Zener voltage regulator ／（2）3-terminal voltage regulator}。
这两件器件正好对应后文的两节：稳压管走 §3--§7，三端稳压器走 §8.4--§8.6。
两者是「同一件事的两种实现」——都是用一只器件把多余的电压吃掉，区别是稳压管\\并联分流、三端稳压器串联降压。

%=============================================================
\hdr{§3}{稳压管：从器件到判据（3.6.1，p15--p18）}

\subhdr{3.1 击穿为什么不等于损坏（p15、p16、p17）}

p15 把两条伏安特性并排：左图是稳压管，第三象限在 $U_{\mathrm{Z}}$ 处\textbf{陡直}下折，
两条水平虚线夹出 $I_{\mathrm{Zmin}}$ 与 $I_{\mathrm{Zmax}}$；右图是普通二极管，
第三象限只有很小的漏电流 $I_{\mathrm{R}}$，一直要到 $U_{\mathrm{BR}}$ 才击穿。
p16 的三段英文说明它的身份：

\callout{\emph{A zener diode is a p-n junction semiconductor device designed to operate in the
reverse breakdown region.}\\
\emph{Zener diode allows electric current in forward direction like a normal diode but also allows
electric current in the reverse direction if the applied reverse voltage is greater than the zener voltage.
Zener diode is always connected in reverse direction because it is specifically designed to work in reverse direction.}}

「\emph{designed to operate}」这四个字是全部关键：\textbf{击穿本身是电学现象，是否损坏是热学问题}。
普通二极管的击穿电流不受控制，结区瞬间耗散功率过大，晶格被烧毁，于是不可逆；
稳压管把击穿电压做得很低、掺杂很重、结面积做得足够大，只要\textbf{限制住电流}（不超过 $I_{\mathrm{Zmax}}$），
击穿就是\textbf{可逆}的——电压撤掉，管子恢复原样。所以「反接」不是接错了，是它的正常工位（p21 的 Notice （1））。

\aside{p17 还给了两条特征 \emph{Reverse breakdown voltage is small}、\emph{Reverse breakdown characteristic is steep}
与五条应用（基准、稳压／分流、开关、削波箝位、保护）。这两条特征不是形容词，各自对应一个可用的量：
「小」对应 $U_{\mathrm{Z}}$ 可以做得很低（几伏），「陡」对应 $r_{\mathrm{Z}}$ 很小（§3.3）。}

\subhdr{3.2 四个参数（p18）：$U_{\mathrm{Z}}$ 是一个范围，不是常数}

p18 用四段文字定义了四个量，本讲义把它们与「怎么用」排在一起：

\par\vspace{2pt}
\noindent\begin{tabular}{@{}p{1.8cm}p{6.4cm}p{7.6cm}@{}}
\tabtop
\thead{参数} & \thead{课件原话（p18）} & \thead{怎么用}
\\
\tabmid
$U_{\mathrm{Z}}$ & \emph{The voltage across the zener diode when it is working}\newline
\emph{It is not a fix value. It's a range.} & 标称值，实际在 $U_{\mathrm{Z}}$ 附近随电流略变；\textbf{不能当成精确常数}
\\
$I_{\mathrm{Zmin}}$ & \emph{If the reverse current is less than this value. There is no regulating effect.}\newline
\emph{Because the reverse voltage is not big enough to force it breakdown.} & \textbf{下限}：电流小于它就没进击穿区，稳压失效
\\
$I_{\mathrm{Zmax}}$ & \emph{If the reverse current is greater than this value. It will be damaged forever.}\newline
\emph{Because the voltage is too big.} & \textbf{上限}：超过它结区过热，永久损坏
\\
$r_{\mathrm{Z}}$ & \emph{Dynamic resistance. The smaller the better.}\newline
$r_{\mathrm{Z}}=\Delta U/\Delta I$ & \textbf{陡度}：越小，同样的电流变化引起的电压变化越小（§3.3）
\tabbot
\end{tabular}

这四段话本身没有错，但\textbf{合起来还缺一个可执行的判据}：给定 $U_{\mathrm{I}}$、$R$、$R_{\mathrm{L}}$，
怎么知道 $I_{\mathrm{Z}}$ 到底落在不在 $I_{\mathrm{Zmin}}$ 与 $I_{\mathrm{Zmax}}$ 之间？
课件没有写这一步，本讲义 §5 补成三步法，§7.2 用它来判 Example 1(g)。

\subhdr{3.3 $r_{\mathrm{Z}}$ 的两条用法：把「陡」变成一个数}

课件只给了定义式与一句「越小越好」，没有说它\textbf{好在哪里}。这里把它用两次，
两次都只用到「在击穿点附近，曲线可线性化」这一条前提。

\textbf{用法一：输入电压变化有多少漏到输出。}在击穿段上把曲线当成切线：
\[
\Delta U_{\mathrm{O}}=r_{\mathrm{Z}}\,\Delta I_{\mathrm{Z}} .
\]
设负载电流不变（$\Delta I_{\mathrm{O}}=0$），由基尔霍夫电流定律 $\Delta I=\Delta I_{\mathrm{Z}}$，
再由回路方程 $\Delta U_{\mathrm{I}}=\Delta I\cdot R+\Delta U_{\mathrm{O}}$，得到
\[
\Delta I_{\mathrm{Z}}=\frac{\Delta U_{\mathrm{I}}}{R+r_{\mathrm{Z}}},
\qquad
\frac{\Delta U_{\mathrm{O}}}{\Delta U_{\mathrm{I}}}=\frac{r_{\mathrm{Z}}}{R+r_{\mathrm{Z}}} .
\]
\textbf{这就是「越小越好」的定量版本}：$r_{\mathrm{Z}}\ll R$ 时，输入怎么波动，输出只按 $r_{\mathrm{Z}}/R$ 的比例跟着动。
代组数：$R=300\,\Omega$、$r_{\mathrm{Z}}=10\,\Omega$ 时比值是 $10/310=3.2\%$；
把 $r_{\mathrm{Z}}$ 换成 $5\,\Omega$，比值降到 $1.6\%$。

\textbf{用法二：负载变化有多少漏到输出。}输入不变（$\Delta U_{\mathrm{I}}=0$），
由 $\Delta I=\Delta I_{\mathrm{Z}}+\Delta I_{\mathrm{O}}$ 与 $\Delta I=-\Delta U_{\mathrm{O}}/R$ 联立，得
\[
\Delta U_{\mathrm{O}}=-(R\parallel r_{\mathrm{Z}})\,\Delta I_{\mathrm{O}} .
\]
即从负载侧看进去，稳压电源的输出电阻是 $R\parallel r_{\mathrm{Z}}$（$R=300\,\Omega$、$r_{\mathrm{Z}}=10\,\Omega$ 时约 $9.7\,\Omega$）。
$r_{\mathrm{Z}}\to0$ 时输出电阻趋于 $0$、$U_{\mathrm{O}}$ 与负载无关——
\textbf{理想稳压管就是把 $r_{\mathrm{Z}}$ 取成零的那只管子}，这也是 §7.3 讲「(f) 用两管相减为什么不算稳压」的伏笔。

\aside{以上两式由 $r_{\mathrm{Z}}=\Delta U/\Delta I$ 与两个基尔霍夫方程推出，$\Delta$ 一律指「相对某个工作点的小变化」。
课件 p18 只写了 $\Delta I=\Delta U/r_{\mathrm{Z}}$ 这一行，且未说明它是\textbf{击穿区切线}的斜率——
按定义它只在击穿段有意义，正向区不能套这个数。
本段与 §4 的两个比值属于本讲义补的推导（[教材补充]），原因：没有它们，
$r_{\mathrm{Z}}$ 只是一个背下来的参数，判不了「这个电源稳不稳」。}

%=============================================================
\hdr{§4}{稳压管稳压电路与稳压过程（p19--p20）}

\subhdr{4.1 电路与三个方程（原形式）}

p19 的电路是本册的骨架：变压器降压、桥式整流、电容 $C$ 滤波，得到 $U_{\mathrm{I}}$；
然后在正母线上串一只限流电阻 $R$，其右侧节点向下接\textbf{反接}的稳压管 $D_{\mathrm{Z}}$，
同一节点向右接负载 $R_{\mathrm{L}}$，$D_{\mathrm{Z}}$ 与 $R_{\mathrm{L}}$ 的下端都回到负母线。
三者的关系是：$C$、$D_{\mathrm{Z}}$、$R_{\mathrm{L}}$ 三者\textbf{并联}在同一对输出端点，$R$ 串在整流器与它们之间。

把它写成本讲义后文一律使用的三个方程：
\[
U_{\mathrm{O}}=U_{\mathrm{I}}-IR\quad(\text{回路电压}),
\qquad
I=I_{\mathrm{Z}}+I_{\mathrm{O}}\quad(\text{节点电流}),
\qquad
U_{\mathrm{O}}=U_{\mathrm{Z}}+r_{\mathrm{Z}}I_{\mathrm{Z}}\quad(\text{击穿段切线}).
\]
第三个方程在 $r_{\mathrm{Z}}\to0$ 时退化为 $U_{\mathrm{O}}=U_{\mathrm{Z}}$，这就是「稳压」二字的全部内容：
\textbf{输出电压等于稳压管的击穿电压，与输入、与负载都无关}。
前两个方程负责回答「那只管子能不能撑住这个工位」。

\aside{课件的 p19、p20 只有图与标注（$U_{\mathrm{I}}$、$U_{\mathrm{O}}$、$I$、$I_{\mathrm{Z}}$、$I_{\mathrm{O}}$），
\textbf{三个方程一条都没有写}。它们是电路的基尔霍夫方程与 p18 的参数定义，本讲义显式写出来，
因为后面所有判断（§5 判据、§6 七个例题、§7 的失败分析）都要直接调用它们。}

\subhdr{4.2 把 p20 的箭头链补成不等式链}

p20 给的是一串箭头（原式照抄，\\ 号表示折行）：

\callout{\emph{Process of voltage regulating：}\\
$U_{\mathrm{O}}\uparrow\ \rightarrow\ I_{\mathrm{Z}}\uparrow\uparrow\ \rightarrow\ I\uparrow\ \rightarrow\ IR\uparrow$
\quad 然后折返回 \quad $U_{\mathrm{O}}\downarrow$}

这五步每一跳都要有出处，否则它只是一句口号。逐跳对上 §4.1 的方程：

\begin{enumerate}
\item $U_{\mathrm{O}}\uparrow\Rightarrow I_{\mathrm{Z}}\uparrow\uparrow$：由 $U_{\mathrm{O}}=U_{\mathrm{Z}}+r_{\mathrm{Z}}I_{\mathrm{Z}}$
得 $\Delta I_{\mathrm{Z}}=\Delta U_{\mathrm{O}}/r_{\mathrm{Z}}$。因为 $r_{\mathrm{Z}}$ 很小，
\textbf{电压的一个小增量引起电流的一个大增量}——p20 用两个箭头「$\uparrow\uparrow$」标的就是这件事，
p18 那句 \emph{A small change of voltage will cause a great change of the current} 是同一句话。
\item $I_{\mathrm{Z}}\uparrow\Rightarrow I\uparrow$：由 $I=I_{\mathrm{Z}}+I_{\mathrm{O}}$，负载电流不变时 $I$ 随 $I_{\mathrm{Z}}$ 同步增。
\item $I\uparrow\Rightarrow IR\uparrow$：$R$ 是常数，直接乘。
\item $IR\uparrow\Rightarrow U_{\mathrm{O}}\downarrow$：由 $U_{\mathrm{O}}=U_{\mathrm{I}}-IR$，$U_{\mathrm{I}}$ 由前级决定、短期不变，
于是 $IR$ 增大多少，$U_{\mathrm{O}}$ 就被压低多少。
\end{enumerate}

四步走完是一个\textbf{负反馈}：输出往上跑一点点，管子立刻多分流一大截，把限流电阻上的压降顶上去，
再把输出压回来。\textbf{稳压管不是「产生」一个电压，而是「吃掉」多余的电压}——
这个说法要在脑子里换过来，否则 (g) 那样的题会算成「管子输出 $10\,\mathrm{V}$，所以不消耗功率」。

\warnbox{\textbf{这条链只说明方向，不说明幅度。}箭头能告诉你「输入升，输出升得少」，
但「少多少」必须用 §3.3 的比值：$\Delta U_{\mathrm{O}}/\Delta U_{\mathrm{I}}=r_{\mathrm{Z}}/(R+r_{\mathrm{Z}})$。
$R=300\,\Omega$、$r_{\mathrm{Z}}=10\,\Omega$ 时这个比值是 $3.2\%$，
即输入涨 $1\,\mathrm{V}$、输出只涨 $32\,\mathrm{mV}$。\textbf{没有这个数，链子只是定性的}。}

\subhdr{4.3 两种扰动，一条链能同时管住}

同一条负反馈链要应付两种完全不同的扰动，这是它值得单独讲的原因：

\textbf{（1）电网电压波动}（$U_{\mathrm{I}}$ 变，$R_{\mathrm{L}}$ 不变）：$U_{\mathrm{I}}\uparrow$ 会让 $I$ 变大，
$I_{\mathrm{Z}}$ 随之增加，多出来的电压被 $R$ 吃掉；输出只按 $r_{\mathrm{Z}}/(R+r_{\mathrm{Z}})$ 的比例跟着动。

\textbf{（2）负载变化}（$U_{\mathrm{I}}$ 不变，$R_{\mathrm{L}}$ 变）：负载变轻（$R_{\mathrm{L}}\uparrow$）时 $I_{\mathrm{O}}$ 减小，
总电流 $I$ 想减小；但稳压管在击穿区「守」着 $U_{\mathrm{O}}$，只要电流还在 $I_{\mathrm{Zmin}}$ 以上，
它就\textbf{把减小的那部分主动接管过来}（$I_{\mathrm{Z}}$ 增大），于是 $I$ 几乎不变、$IR$ 几乎不变、$U_{\mathrm{O}}$ 几乎不变。

第（2）种情况有一个必须记住的极限：$I_{\mathrm{Z}}$ 不能无限增大。
负载轻到端子上没有电流时，全部 $I$ 都从稳压管走，此时
\[
I_{\mathrm{Z,max}}\big|_{\text{空载}}=\frac{U_{\mathrm{I}}-U_{\mathrm{Z}}}{R}.
\]
若这个数超过管子的 $I_{\mathrm{Zmax}}$，\textbf{空载时管子会被烧}（而不是像 §1.2 的电容那样只是电压上升）。
所以「限流电阻 $R$ 该取多大」这个问题，答案由\textbf{两个最坏工况}夹出来，这就是 §5.2。

%=============================================================
\hdr{§5}{工作点判据与限流电阻的取法}

\subhdr{5.1 三步法：拿到任何稳压管电路都这么判}

§3.2 的四个参数、§4.1 的三个方程，合起来给出一个三步判据。
它的思路是\textbf{先假设稳压成立，再验算这个假设能不能站住}：

\par\vspace{2pt}
\noindent\begin{tabular}{@{}p{1.1cm}p{14.6cm}@{}}
\tabtop
\thead{步} & \thead{动作}
\\
\tabmid
1 & \textbf{假定稳压成立}：$U_{\mathrm{O}}=U_{\mathrm{Z}}$。算出总电流 $I=(U_{\mathrm{I}}-U_{\mathrm{Z}})/R$
与负载电流 $I_{\mathrm{O}}=U_{\mathrm{Z}}/R_{\mathrm{L}}$。
\\
2 & \textbf{求稳压管电流}：$I_{\mathrm{Z}}=I-I_{\mathrm{O}}$。
\\
3 & \textbf{验算三选一}：
$I_{\mathrm{Zmin}}\le I_{\mathrm{Z}}\le I_{\mathrm{Zmax}}$ 则假定成立，$U_{\mathrm{O}}=U_{\mathrm{Z}}$；
$I_{\mathrm{Z}}>I_{\mathrm{Zmax}}$ 则管子过热损坏（$U_{\mathrm{O}}$ 塌）；
$I_{\mathrm{Z}}<I_{\mathrm{Zmin}}$（含 $I_{\mathrm{Z}}\le0$）则根本没进击穿区，管子相当于断开，$U_{\mathrm{O}}\approx U_{\mathrm{I}}$。
\tabbot
\end{tabular}

第 3 步的第三种情形值得单独说：\textbf{一旦退出击穿，$R$ 上几乎没有电流流过，也就几乎没有压降}，
所以输出不是「比 $U_{\mathrm{Z}}$ 小一点」，而是「几乎等于 $U_{\mathrm{I}}$」——
这是本题组最常见的错解来源（§6.9 与 §7.1 都回到这一点）。

\aside{三步法的第 1 步与第 3 步的顺序不能颠倒：代入 $U_{\mathrm{O}}=U_{\mathrm{Z}}$ 之前，
电路里 $I_{\mathrm{Z}}$ 是未知量，无法「先算 $I_{\mathrm{Z}}$ 再算 $U_{\mathrm{O}}$」。
课件 p18 把 $I_{\mathrm{Zmin}}$、$I_{\mathrm{Zmax}}$ 写成两句条件，没有给这个「假设--验算」的结构。}

\subhdr{5.2 限流电阻 $R$ 的窗口：两条最坏工况}

设计问题通常是这样问的：输入在 $U_{\mathrm{Imin}}\sim U_{\mathrm{Imax}}$ 之间波动，负载电流在
$0\sim I_{\mathrm{Omax}}$ 之间变化，要求稳压管始终工作在 $I_{\mathrm{Zmin}}\sim I_{\mathrm{Zmax}}$ 内，
求 $R$ 的取值范围。

\textbf{上界（最轻载 + 最低输入）}：此时 $I_{\mathrm{Z}}$ 最小，要保证它不低于下限
\[
\frac{U_{\mathrm{Imin}}-U_{\mathrm{Z}}}{R}\ge I_{\mathrm{Zmin}}+I_{\mathrm{Omax}}
\ \Longrightarrow\
R\le\frac{U_{\mathrm{Imin}}-U_{\mathrm{Z}}}{I_{\mathrm{Zmin}}+I_{\mathrm{Omax}}} .
\]

\textbf{下界（空载 + 最高输入）}：此时 $I_{\mathrm{Z}}$ 最大，要保证它不超过上限
\[
\frac{U_{\mathrm{Imax}}-U_{\mathrm{Z}}}{R}\le I_{\mathrm{Zmax}}+I_{\mathrm{Omin}}
\ \Longrightarrow\
R\ge\frac{U_{\mathrm{Imax}}-U_{\mathrm{Z}}}{I_{\mathrm{Zmax}}+I_{\mathrm{Omin}}} .
\]

两条一夹，$R$ 只能落在一个\textbf{区间}里。区间为空，说明这只稳压管配这个输入范围根本做不出来，
要么换管、要么改输入——这是设计题真正的答案形态，不是「一个数」。

\subhdr{5.3 一个完整算例（自编，非往年题）}

设 $U_{\mathrm{Z}}=10\,\mathrm{V}$、$I_{\mathrm{Zmin}}=5\,\mathrm{mA}$、$I_{\mathrm{Zmax}}=23\,\mathrm{mA}$，
输入在 $18\sim22\,\mathrm{V}$ 之间波动（电网 $\pm10\%$ 的常见量级），负载电流在 $0\sim6.67\,\mathrm{mA}$ 之间变化
（对应 $R_{\mathrm{L}}\ge1.5\,\mathrm{k\Omega}$）。求 $R$。

代上面两条：
\[
R\ge\frac{22-10}{0.023+0}=521.7\,\Omega,
\qquad
R\le\frac{18-10}{0.005+0.00667}=685.7\,\Omega .
\]
所以 $R\in[522,\,686]\,\Omega$，取标称 $620\,\Omega$。验算两端：

\begin{itemize}
\item 最高输入、空载：$I_{\mathrm{Z}}=(22-10)/620=19.4\,\mathrm{mA}\le23\,\mathrm{mA}$，合格；
\item 最低输入、满载：$I_{\mathrm{Z}}=(18-10)/620-6.67=6.24\,\mathrm{mA}\ge5\,\mathrm{mA}$，合格。
\end{itemize}

再补一步\textbf{功率校验}（很多题会漏）：$P_{R}=(22-10)^2/620=0.23\,\mathrm{W}$，取 $1/2\,\mathrm{W}$ 电阻；
$P_{Z}=U_{\mathrm{Z}}I_{\mathrm{Zmax}}=10\times0.023=0.23\,\mathrm{W}$，所选管子的额定功耗必须大于它。
\aside{本例的四个输入（$U_{\mathrm{Z}}$、$I_{\mathrm{Zmin}}$、$I_{\mathrm{Zmax}}$、输入范围）取自 p18 的参数含义与 p27 的题面量级，
$18\sim22\,\mathrm{V}$ 的波动范围是本讲义设定的。此例用于演示 §5.2 的窗口怎么用，不是课件例题。}

\subhdr{5.4 三个常见错法}

\begin{enumerate}
\item \textbf{把 $R$ 当自由度随便取。} $R$ 小了烧管（空载时 $I_{\mathrm{Z}}$ 超限），$R$ 大了失稳
（满载时 $I_{\mathrm{Z}}$ 掉到 $I_{\mathrm{Zmin}}$ 以下）。它是被夹出来的，不是选的。
\item \textbf{只算一种工况。} 只按额定输入算 $R$，电网一波动就出窗口——必须用\textbf{两个最坏值}各算一次。
\item \textbf{忘了功率。} $R$ 与稳压管都要按最坏工况算耗散，否则电路「算对了」但元件「受不了」。
\end{enumerate}

%=============================================================
\hdr{§6}{主例题 Example 1 全程（p21--p27）}

\subhdr{6.1 题面结构：七个小问在回答 p21 的两个 Notice}

七个电路的公共结构是：输入 $U_{\mathrm{I}}$ 经 $R$ 接到节点，节点向下串两只稳压管 $D_{\mathrm{Z1}}$、$D_{\mathrm{Z2}}$
（\textbf{都反接}，即阴极朝正母线），输出 $U_{\mathrm{O}}$ 取自这一串的两端。已知量全组相同：
\[
U_{\mathrm{Z1}}=4.5\,\mathrm{V},\qquad U_{\mathrm{Z2}}=9.5\,\mathrm{V},\qquad
U_{\mathrm{F}}=0.5\,\mathrm{V}\ (\text{正向压降}).
\]
所以两只管子\textbf{反接串联}时稳压值是 $4.5+9.5=14\,\mathrm{V}$，\textbf{正接串联}时压降是 $0.5+0.5=1\,\mathrm{V}$。

p21 页脚提出了两个 Notice：\emph{（1）In order to stabilize the voltage, the zener diode need to be connected in reverse.}
与 \emph{（2）When will it fail to stabilize the voltage?} 七个小问就是这两个提问的五种答案：

\par\vspace{2pt}
\noindent\begin{tabular}{@{}p{1.4cm}p{2.2cm}p{4.5cm}p{7.3cm}@{}}
\tabtop
\thead{问} & \thead{页码} & \thead{已知} & \thead{它在演哪一种情形}
\\
\tabmid
(a) & p21 & $U_{\mathrm{I}}=18\,\mathrm{V}$，正接 & 正常稳压：$18>14$，两管都在击穿区
\\
(b) & p22 & $U_{\mathrm{I}}=10\,\mathrm{V}$，正接 & \textbf{输入不足}：$10<14$，进不了击穿区
\\
(c) & p23 & $U_{\mathrm{I}}=4\,\mathrm{V}$，输入反接 & 反接 = 正向导通，两只各掉 $0.5\,\mathrm{V}$
\\
(d) & p24 & $U_{\mathrm{I}}=4\,\mathrm{V}$，输入反接 & 同 (c)（两页位图标注相同，差异标 [待确认]）
\\
(e) & p25 & $U_{\mathrm{I}}=8\,\mathrm{V}$，正接 & 又一种输入不足：$8<14$
\\
(f) & p26 & $U_{\mathrm{I}}=13\,\mathrm{V}$，两支路 & 两节点相减：$4.5$ 与 $9.5$ 之差
\\
(g) & p27 & $U_{\mathrm{I}}=20\,\mathrm{V}$，$R=300\,\Omega$ & \textbf{带载}：判 $I_{\mathrm{Z}}$ 是否超限
\tabbot
\end{tabular}

\subhdr{6.2 （a）$U_{\mathrm{I}}=18\,\mathrm{V}$：正常稳压，输出 $14\,\mathrm{V}$}

反接串联的两只管子要同时击穿，需要 $\ge U_{\mathrm{Z1}}+U_{\mathrm{Z2}}=14\,\mathrm{V}$。
现在 $U_{\mathrm{I}}=18\,\mathrm{V}>14\,\mathrm{V}$，两管都进击穿区，输出被钳在
\[
U_{\mathrm{O}}=U_{\mathrm{Z1}}+U_{\mathrm{Z2}}=4.5+9.5=14\,\mathrm{V}.
\]
分压怎么分？每一只都在自己的击穿点上，所以 $D_{\mathrm{Z1}}$ 上 $4.5\,\mathrm{V}$、$D_{\mathrm{Z2}}$ 上 $9.5\,\mathrm{V}$，
\textbf{不是按电阻比分的}——这是「串联稳压值相加」的物理含义：击穿电压由器件决定，不由外电路决定。

剩下 $18-14=4\,\mathrm{V}$ 落在 $R$ 上，故 $I=(18-14)/R$。这个电流全部流过稳压管（本页没有接负载），
所以是否安全取决于 $R$：课件未给 $R$ 与 $I_{\mathrm{Zmax}}$，无法验算，标 [待确认]。
\aside{若取 $R=100\,\Omega$，则 $I_{\mathrm{Z}}=40\,\mathrm{mA}$。对照 p27 给出的量级（$I_{\mathrm{Zmax}}=23\,\mathrm{mA}$），
这个 $R$ 偏小——可见 (a) 只问「输出是多少」，没打算让你验算电流。本注是补充核算，不是例题要求。}

\subhdr{6.3 （b）$U_{\mathrm{I}}=10\,\mathrm{V}$：进不了击穿区，输出等于输入}

$10\,\mathrm{V}<14\,\mathrm{V}$，两只管子都停在击穿之前，只有极小的反向漏电流。
漏电流在 $R$ 上的压降可以忽略，于是
\[
U_{\mathrm{O}}\approx U_{\mathrm{I}}=10\,\mathrm{V}.
\]
\textbf{注意这里不是「$14\,\mathrm{V}$ 减去什么」}，而是整个支路相当于断开，输入原样出现在输出端。
这就是 Notice（2）的第一个答案：\textbf{输入电压不够高时，稳压失效，但输出不会消失，而是跟着输入跑。}

\aside{严格的写法是 $U_{\mathrm{O}}=U_{\mathrm{I}}-IR$，而 $I=I_{\mathrm{Z}}+I_{\mathrm{O}}$ 中的 $I_{\mathrm{Z}}$ 只有漏电流（$\mu\mathrm{A}$ 量级），
$I_{\mathrm{O}}=0$（本页无负载），所以 $IR\approx0$，$U_{\mathrm{O}}\approx U_{\mathrm{I}}$。
课件没有写这一步，只给了电路；把它补出来，「为什么是 $10\,\mathrm{V}$ 而不是别的数」才说得通。}

\subhdr{6.4 （c）（d）$U_{\mathrm{I}}=4\,\mathrm{V}$，输入反接：两只管子正向导通}

p23 与 p24 的端子标注与前两页相反：\textbf{输入端上「$-$」下「$+$」}（输出端仍标上「$+$」下「$-$」）。
这一反接改变了整条支路的工作方式：电流从底端（$+$）出发，经 $D_{\mathrm{Z2}}$、$D_{\mathrm{Z1}}$ 的\textbf{正向}
流向顶端，两只管子都变成普通正向导通的二极管。

正向导通后每只压降 $0.5\,\mathrm{V}$，串联共 $1\,\mathrm{V}$，于是底端比顶端高 $1\,\mathrm{V}$；
按图上标注的参考极性（顶端为 $+$、底端为 $-$），
\[
U_{\mathrm{O}}=-(U_{\mathrm{F1}}+U_{\mathrm{F2}})=-1\,\mathrm{V},
\qquad |U_{\mathrm{O}}|=1\,\mathrm{V}.
\]
这是 Notice（2）的第二个答案：\textbf{接反了就不稳压}——不但不稳压，输出电压还会反号、只剩 $1\,\mathrm{V}$。

\warnbox{\textbf{$4\,\mathrm{V}$ 在这里几乎不起作用。}只要 $U_{\mathrm{I}}>1\,\mathrm{V}$，两只正向导通的管子就把压降钉在 $1\,\mathrm{V}$，
多余的部分全落在 $R$ 上，即 $U_{R}=4-1=3\,\mathrm{V}$。
换句话说：(c)(d) 的输出\textbf{与 $U_{\mathrm{I}}$ 无关}，这是「正向钳位」与「反向稳压」两种机制都有的性质，
但钳位值是 $1\,\mathrm{V}$ 而不是 $14\,\mathrm{V}$。}

\aside{p23 与 p24 两页的位图标注（输入上「$-$」下「$+$」、输出上「$+$」下「$-$」）逐字相同，
扫描描述也只差一处：p24 把「$+$」标在 $R$ 之后、稳压管顶端那个节点上。
两者差异在本册的文字层与扫描图里都无法确定，标 [待确认]：本讲义按「两页都是输入反接」处理，
答案同为 $|U_{\mathrm{O}}|=1\,\mathrm{V}$。若考试遇到，先看\textbf{输入端}的极性标记，那是唯一决定正/反向的量。}

\subhdr{6.5 （e）$U_{\mathrm{I}}=8\,\mathrm{V}$：又一种输入不足}

正接（输入端上「$+$」下「$-$」），$8\,\mathrm{V}<14\,\mathrm{V}$，情形与 (b) 完全相同：
\[
U_{\mathrm{O}}\approx U_{\mathrm{I}}=8\,\mathrm{V}.
\]
(b) 与 (e) 摆在一起说明一件事：\textbf{「不稳压」不是一个开关，而是一段区间}——
正接时只要 $0<U_{\mathrm{I}}<14\,\mathrm{V}$，输出就从 $U_{\mathrm{I}}$ 一路跟着输入爬，
直到 $U_{\mathrm{I}}\ge14\,\mathrm{V}$ 且 $I_{\mathrm{Z}}$ 不超限，才真正被钉在 $14\,\mathrm{V}$ 上。
§6.4 那种反接是另一回事：它把输出钉在 $-1\,\mathrm{V}$，与 $U_{\mathrm{I}}$ 的大小无关。

\subhdr{6.6 （f）$U_{\mathrm{I}}=13\,\mathrm{V}$，两支路：输出取两节点之差，$5\,\mathrm{V}$}

p26 把两只管子拆成\textbf{两条并列支路}：左支路 $R+D_{\mathrm{Z1}}$、右支路 $R+D_{\mathrm{Z2}}$，
输出 $U_{\mathrm{O}}$ 取两条支路\textbf{中间两个节点}之间的电压（左节点标 $+$、右节点标 $-$）。
两条支路各自与 §4 的电路完全同构，各按自己的稳压值工作：

\begin{itemize}
\item 左支路：$13>4.5$，$D_{\mathrm{Z1}}$ 击穿，左节点电位 $U_{\mathrm{L}}=4.5\,\mathrm{V}$；
\item 右支路：$13>9.5$，$D_{\mathrm{Z2}}$ 击穿，右节点电位 $U_{\mathrm{R}}=9.5\,\mathrm{V}$。
\end{itemize}

按标注的参考极性，
\[
U_{\mathrm{O}}=U_{\mathrm{L}}-U_{\mathrm{R}}=4.5-9.5=-5\,\mathrm{V},
\qquad |U_{\mathrm{O}}|=5\,\mathrm{V}.
\]
机制是\textbf{两个稳压值相减}：$9.5-4.5=5$。它说明稳压值可以「凑」（用一只高价管减一只低价管），
代价是\textbf{两个节点的绝对精度都会带进差值}——这也是 §7.3 要展开的。

\aside{两支路各要满足 $I_{\mathrm{Zmin}}\le I_{\mathrm{Z}}\le I_{\mathrm{Zmax}}$，判据同 §5.1。
课件未给 $R$、$I_{\mathrm{Zmin}}$、$I_{\mathrm{Zmax}}$，所以「两条支路都真的稳压」这一步无法验算，标 [待确认]。
若某支路电流不够，该节点电位就不再是 $U_{\mathrm{Z}}$ 而是 $U_{\mathrm{I}}$，差值立刻变成另一个数——
这正是 (f) 这类电路在实际中不受欢迎的原因。}

\subhdr{6.7 （g）$U_{\mathrm{I}}=20\,\mathrm{V}$，$R=300\,\Omega$，$R_{\mathrm{L}}=1.5\,\mathrm{k\Omega}$：判 $I_{\mathrm{Z}}$}

p27 的题面是 \emph{$U_{\mathrm{Z}}=10\,\mathrm{V}$,\ $I_{\mathrm{Zmax}}=23\,\mathrm{mA}$. Find $I_{\mathrm{Z}}$}，
电路是 $R=300\,\Omega$ 串在输入、$D_{\mathrm{Z}}$ 与 $R_{\mathrm{L}}=1.5\,\mathrm{k\Omega}$ 并联。按 §5.1 三步法：

\textbf{第 1 步}，假定稳压成立，$U_{\mathrm{O}}=U_{\mathrm{Z}}=10\,\mathrm{V}$：
\[
I=\frac{U_{\mathrm{I}}-U_{\mathrm{Z}}}{R}=\frac{20-10}{300}=33.3\,\mathrm{mA},
\qquad
I_{\mathrm{O}}=\frac{U_{\mathrm{Z}}}{R_{\mathrm{L}}}=\frac{10}{1500}=6.67\,\mathrm{mA}.
\]
\textbf{第 2 步}，节点电流：
\[
I_{\mathrm{Z}}=I-I_{\mathrm{O}}=33.3-6.67=26.7\,\mathrm{mA}.
\]
\textbf{第 3 步}，验算：$I_{\mathrm{Z}}=26.7\,\mathrm{mA}>I_{\mathrm{Zmax}}=23\,\mathrm{mA}$。
\textbf{假定不成立}——管子超过最大工作电流，会过热损坏。

所以本题的正确答案不是一个电流值，而是两句：（1）若强行稳压，管子要过 $26.7\,\mathrm{mA}$，超标；
（2）管子一旦烧毁，通常\textbf{短路}，于是 $U_{\mathrm{O}}=0$，$R$ 上单独承受全部 $20\,\mathrm{V}$、
耗散 $20^2/300=1.33\,\mathrm{W}$，很快也会跟着出问题。

要让它正常工作，必须加大 $R$。按最坏工况反解（负载始终保持 $6.67\,\mathrm{mA}$）：
\[
R\ge\frac{20-10}{0.023+0.00667}=337\,\Omega .
\]
取标称 $390\,\Omega$ 验算：$I=(20-10)/390=25.6\,\mathrm{mA}$，$I_{\mathrm{Z}}=25.6-6.67=19.0\,\mathrm{mA}<23\,\mathrm{mA}$，合格；
此时 $P_{R}=10^2/390=0.26\,\mathrm{W}$、$P_{Z}=10\times0.019=0.19\,\mathrm{W}$，都留有余量。

\aside{「管子烧毁后通常短路」不是定理，而是这类器件的常见失效模式（结区熔融形成低阻通路）。
本讲义的 $1.33\,\mathrm{W}$ 与 $390\,\Omega$ 的改法都是补充核算，课件 p27 只问了 $I_{\mathrm{Z}}$ 一个量。}

\subhdr{6.8 七问答案表}

\par\vspace{2pt}
\noindent\begin{tabular}{@{}p{1.2cm}p{3.3cm}p{2.1cm}p{2.9cm}p{5.8cm}@{}}
\tabtop
\thead{问} & \thead{输入（极性）} & \thead{状态} & \thead{$U_{\mathrm{O}}$} & \thead{依据}
\\
\tabmid
(a) & $18\,\mathrm{V}$（正接） & 两管击穿 & $14\,\mathrm{V}$ & $U_{\mathrm{Z1}}+U_{\mathrm{Z2}}$（p21）
\\
(b) & $10\,\mathrm{V}$（正接） & 未击穿 & $10\,\mathrm{V}$ & $U_{\mathrm{O}}\approx U_{\mathrm{I}}$（p22）
\\
(c) & $4\,\mathrm{V}$（反接） & 两管正向 & $-1\,\mathrm{V}$ & $-(U_{\mathrm{F1}}+U_{\mathrm{F2}})$（p23）
\\
(d) & $4\,\mathrm{V}$（反接） & 两管正向 & $-1\,\mathrm{V}$ & 同 (c)，[待确认]（p24）
\\
(e) & $8\,\mathrm{V}$（正接） & 未击穿 & $8\,\mathrm{V}$ & $U_{\mathrm{O}}\approx U_{\mathrm{I}}$（p25）
\\
(f) & $13\,\mathrm{V}$（两支路） & 各自击穿 & $-5\,\mathrm{V}$ & $4.5-9.5$（p26）
\\
(g) & $20\,\mathrm{V}$（带载） & $I_{\mathrm{Z}}$ 超限 & 不能稳压 & $26.7\,\mathrm{mA}>23\,\mathrm{mA}$（p27）
\tabbot
\end{tabular}

\subhdr{6.9 错解诊断：每个错法会得到什么数}

\par\vspace{2pt}
\noindent\begin{tabular}{@{}p{5.3cm}p{3.1cm}p{2.7cm}p{4.3cm}@{}}
\tabtop
\thead{错法} & \thead{会算出} & \thead{正确} & \thead{错在哪一步}
\\
\tabmid
(b) 用「$14-U_{\mathrm{I}}$」或「$U_{\mathrm{I}}/14$」 & $4\,\mathrm{V}$ 之类 & $10\,\mathrm{V}$ & 把支路当成仍在击穿区
\\
(c)(d) 认为反接也能稳压 & $14\,\mathrm{V}$ & $-1\,\mathrm{V}$ & 忘了反接意味着正向导通
\\
(e) 认为「不足就输出 $0$」 & $0\,\mathrm{V}$ & $8\,\mathrm{V}$ & 把「失效」当成「断路」
\\
(f) 用 $4.5+9.5$ & $14\,\mathrm{V}$ & $-5\,\mathrm{V}$ & 没有看输出取的是\textbf{两节点之差}
\\
(g) 直接写 $I_{\mathrm{Z}}=26.7\,\mathrm{mA}$ 就交卷 & $26.7\,\mathrm{mA}$ & 判超限 & 漏了第 3 步验算
\tabbot
\end{tabular}

五个错里有三个同一个病根：\textbf{把「稳压管支路」当成一个固定电压源}。
它只在击穿区里是电压源，出了击穿区要么是断开的（(b)(e)）、要么是 $1\,\mathrm{V}$ 的正向压降（(c)(d)）。
§7 把这三处的「算不下去」讲透。

%=============================================================
\hdr{§7}{三处「算不下去」的地方}

本节回答三个「课件给了图、没给为什么」的问题：(b)(e) 为什么只能给约等号、(g) 为什么答案不是电流、
(f) 为什么得到的 $5\,\mathrm{V}$ 不能当稳压电源用。

\subhdr{7.1 (b)(e)：方程在这一区无解，只能给不等式}

在 (b) 里，如果仍按 §4.1 的第三个方程 $U_{\mathrm{O}}=U_{\mathrm{Z}}$ 算，会立刻矛盾：
$U_{\mathrm{O}}=14\,\mathrm{V}$ 而 $U_{\mathrm{I}}=10\,\mathrm{V}$，串联支路上「输入小于输出」，电流只能为负。
这说明\textbf{假设本身出错了}，而不是算术出错。

真实情况是：两只管子都在击穿之前，伏安特性是一条几乎水平的线，等效成一只\textbf{阻值极大}的电阻。
此时电路是两个大阻值电阻（$R$ 与管子）串联分压，$R$ 上的压降是 $\mu\mathrm{A}$ 量级的漏电流乘 $R$，可以忽略，
于是 $U_{\mathrm{O}}\approx U_{\mathrm{I}}$。

\textbf{为什么只能在约等号上停下}：要算出精确值，需要管子在击穿前的\textbf{漏电阻}（或反向饱和电流），
这不是稳压管的标称参数，课件的四个参数里也没有它。因此这道题的\textbf{答案形态}就是
「$U_{\mathrm{O}}\approx U_{\mathrm{I}}$，稳压失效」，再往下算只是在编数据。
这就是缺口清单第六类（答案形态）的现场：\textbf{停在这里不是偷懒，是数据不够}。

\subhdr{7.2 (g)：题目问的是电流，正确答案是「会坏」}

p27 的题干是 \emph{Find $I_{\mathrm{Z}}$}，字面上只求一个数。但三步法走完会发现
$I_{\mathrm{Z}}=26.7\,\mathrm{mA}$ 与 $I_{\mathrm{Zmax}}=23\,\mathrm{mA}$ 的大小关系才是本题的落点。
\textbf{题目给 $I_{\mathrm{Zmax}}$ 这个条件是本题唯一的「考点提示」}：
凡是题面同时给出 $U_{\mathrm{Z}}$、$I_{\mathrm{Zmax}}$、$R$、$R_{\mathrm{L}}$ 与 $U_{\mathrm{I}}$，问「求 $I_{\mathrm{Z}}$」，
十有八九是要你判超限（或判不足），而不是要你报一个电流值。

这与 §5.1 第 3 步的三种情形一一对应：
\[
I_{\mathrm{Z}}>I_{\mathrm{Zmax}}\Rightarrow\text{损坏};\qquad
I_{\mathrm{Z}}<I_{\mathrm{Zmin}}\Rightarrow\text{失稳};\qquad
I_{\mathrm{Zmin}}\le I_{\mathrm{Z}}\le I_{\mathrm{Zmax}}\Rightarrow\text{稳压成立}.
\]
本题落在第一档，所以答题必须写「不稳压，管子过流」，再补一句「需增大 $R$ 至 $\ge337\,\Omega$」才算完整——
\textbf{只写 $26.7\,\mathrm{mA}$ 是把一道判断题做成了计算题}。

\subhdr{7.3 (f)：$5\,\mathrm{V}$ 为什么不能当稳压电源用}

(f) 用两只管子相减凑出 $5\,\mathrm{V}$，算术上成立，工程上有三处硬伤：

\textbf{（1）误差相加，不是相减。} $U_{\mathrm{Z}}$ 本身是一个范围（p18 原话 \emph{It is not a fix value.
It's a range.}）。若两只管子的偏差各为 $\pm5\%$，
\[
4.5\pm0.225\,\mathrm{V},\qquad 9.5\pm0.475\,\mathrm{V},
\]
最坏情况下差值是 $5\pm0.7\,\mathrm{V}$，\textbf{相对误差 $\pm14\%$}。
差值把两处误差\textbf{叠加}了，比单用一只 $5\,\mathrm{V}$ 的管子差得多（后者只有 $\pm5\%$）。

\textbf{（2）两条支路各自都要满足电流条件。} $U_{\mathrm{I}}=13\,\mathrm{V}$ 对 $9.5\,\mathrm{V}$ 那条支路只剩 $3.5\,\mathrm{V}$ 落在 $R$ 上，
对 $4.5\,\mathrm{V}$ 那条剩 $8.5\,\mathrm{V}$。同一个 $R$ 下两条支路的电流差近三倍，
很可能一条超 $I_{\mathrm{Zmax}}$、另一条低于 $I_{\mathrm{Zmin}}$——此时两个节点电位都不再等于 $U_{\mathrm{Z}}$，差值立刻失效。
课件未给 $R$，无法验算，标 [待确认]。

\textbf{（3）输出「浮」在两个非地节点之间。} $U_{\mathrm{O}}$ 既不接正母线也不接地，
后级电路若要共地就得另做处理——这是实际电路里最不受欢迎的一种输出方式。

三处合起来给出一个可迁移的判断：\textbf{稳压值可以相减，但代价是精度与电流预算两项都要重新算}。
与其相减，不如直接选一只 $5\,\mathrm{V}$ 的稳压管。

%=============================================================
\hdr{§8}{Example 2 与三端集成稳压器（p28--p38）}

\subhdr{8.1 p28--p36：同一个例题的九帧}

p28 起小标题写的是 \emph{4. Example 2}，电路与前面完全不同：\textbf{变压器 + 桥式整流 + 电容滤波 + 限流电阻 $R$ +
稳压管 + 负载 $R_{\mathrm{L}}$} 的整机图，右侧一只数字万用表，红表笔接输出正端、黑表笔接负端，\textbf{读数 $000$}。
p29--p36 是同一张底图的连续帧，只有三点变化：

\begin{enumerate}
\item 手写笔画的颜色与位置（桥臂上一道洋红、电容支路一道红、稳压管支路一道红、负载支路一道浅蓝）——
说明老师在讲台上逐段指过这条链；
\item 页眉的章节号：p28 写 \emph{3.6}，p29 写 \emph{8.6.1}，p30--p36 写 \emph{8.6 Voltage regulating circuit}；
\item p28 还并列着那张四列总表（同 §2.1）。
\end{enumerate}

\warnbox{\textbf{章节号在册内跳了}：同一节内容，p28 记作 3.6，p29 起改记 8.6。这是课件的编号笔误，
按内容应统一为 3.6（与 p15--p27 一致）。本讲义引用时一律写成 3.6，
涉及该处只标页码，不沿用 8.6 的编号。}

九帧里\textbf{没有给任何数}（$U_2$、$U_{\mathrm{Z}}$、$R$、$R_{\mathrm{L}}$ 全无标注），也没有解题过程。
所以本节做两件课件没做的事：把这条整机链的关系写成一套可代入的式子（§8.2），
再解释「读数 $000$」在什么情况下才是真实故障（§8.3）。

\subhdr{8.2 整机链的设计关系：从 $U_2$ 一路算到 $I_{\mathrm{Z}}$}

p28 的电路可以按「两级」写成一条链，前一级用 §2 的表，后一级用 §5 的判据：

\par\vspace{2pt}
\noindent\begin{tabular}{@{}p{1.0cm}p{3.6cm}p{11.1cm}@{}}
\tabtop
\thead{级} & \thead{式子} & \thead{说明}
\\
\tabmid
（1） & $U_{\mathrm{I}}=1.2U_2$ & 桥式 + 电容滤波（§2.1 表，全波有电容一列）
\\
（2） & $U_{\mathrm{O}}=U_{\mathrm{Z}}$ & 稳压管在击穿区工作（§5.1 第 1 步，须事后验算）
\\
（3） & $I=(U_{\mathrm{I}}-U_{\mathrm{Z}})/R$ & 限流电阻上的电流
\\
（4） & $I_{\mathrm{O}}=U_{\mathrm{Z}}/R_{\mathrm{L}}$ & 负载电流
\\
（5） & $I_{\mathrm{Z}}=I-I_{\mathrm{O}}$ & 负反馈链（§4.2）的落点，须满足 $I_{\mathrm{Zmin}}\le I_{\mathrm{Z}}\le I_{\mathrm{Zmax}}$
\tabbot
\end{tabular}

这五步是全册的「自检链」：任何一道整机题，只要顺着走一遍，就能同时得到输出电压与元件应力。
反方向也成立——若题给 $U_{\mathrm{O}}$、$I_{\mathrm{O}}$ 与输入范围，就是 §5.2 的 $R$ 窗口问题。

\subhdr{8.3 读数 $000$：什么时候是真故障}

万用表读 $000$ 有两种完全不同的含义，必须先分清，否则会误判：
\textbf{（甲）表没工作}（挡位错、表笔没接触、电池没电）——此时 $000$ 不含任何电路信息；
\textbf{（乙）输出确实为零}——真故障。判断方法：先量 $U_{\mathrm{I}}$（电容两端），再量 $U_{\mathrm{O}}$，两步就分开了。

\par\vspace{2pt}
\noindent\begin{tabular}{@{}p{4.6cm}p{3.4cm}p{7.7cm}@{}}
\tabtop
\thead{测得} & \thead{判断} & \thead{可能原因}
\\
\tabmid
$U_{\mathrm{I}}=0$ & 前级全无电 & 变压器一次侧断电、保险丝断、整流桥两臂同时开路
\\
$U_{\mathrm{I}}\approx0.9U_2$ 且纹波大 & 滤波失效 & 电容开路（对应 §2.4 第（4）问：跌到 $22.5\,\mathrm{V}$ 那种情形）
\\
$U_{\mathrm{I}}$ 正常、$U_{\mathrm{O}}=0$ & 后级短路或断路 & \textbf{稳压管击穿短路}（最常见，见 §6.7）、限流电阻 $R$ 开路、负载短路
\\
$U_{\mathrm{I}}$ 正常、$U_{\mathrm{O}}\approx U_{\mathrm{I}}$ & 稳压失效 & 输入低于 $U_{\mathrm{Z}}$（§6.3 那种情形）、稳压管开路
\tabbot
\end{tabular}

最后一行值得记：\textbf{「输出等于输入」是稳压失效的典型指纹}，
它的成因不是「管子坏了」而是「管子没进击穿区」——这与「输出为零」是两件完全不同的事故。

\aside{课件 p28--p36 只给电路与读数，没有故障分析与判断表。本节的关系链与诊断表是本讲义补的
（[教材补充]），补的原因：一张带万用表读数的图，若不给「怎么读这个数」，它就只是一张装饰画。
表中「$U_{\mathrm{I}}\approx0.9U_2$」一行的数值依据是 §2.1 的表（全波无电容）。}

\subhdr{8.4 集成稳压器的分类（p37）}

p37 是一张树状分类图，标题 \emph{（1）Introduction}\ /\ \emph{Classification}：

\par\vspace{2pt}
\noindent\begin{tabular}{@{}p{5.6cm}p{5.4cm}p{4.7cm}@{}}
\tabtop
\thead{大类} & \thead{子类} & \thead{典型器件}
\\
\tabmid
三端\textbf{固定}稳压器 & 正输出 78XX & 7805、7812、7815
\\
（output 不可调） & 负输出 79XX & 7905、7912、7915
\\
三端\textbf{可调}稳压器 & 正输出 & LM317
\\
（output 可调） & 负输出 & LM337
\tabbot
\end{tabular}

分类的两个维度是「输出可不可调」与「输出正还是负」。课件只画了框图，
器件型号一列是本讲义补的（[教材补充]），补的原因：p38 的实物照片上就印着 7805、7905，
不给型号这一列，分类图与照片对不上号。

\subhdr{8.5 三端稳压器的管脚（p38）：78 与 79 顺序不同}

p38 的标题是 \emph{（2）Three-terminal voltage regulator}\ /\ \emph{（1）Shape and pins}，
左右并列两组：外形图（三根引脚标 $1$、$2$、$3$）加实物照片（7805 与 7905），
两组下方的管脚定义是\textbf{不一样}的：

\par\vspace{2pt}
\noindent\begin{tabular}{@{}p{2.6cm}p{2.7cm}p{2.7cm}p{2.3cm}p{4.9cm}@{}}
\tabtop
\thead{系列} & \thead{1 脚} & \thead{2 脚} & \thead{3 脚} & \thead{输出极性}
\\
\tabmid
78XX & IN & GND & OUT & 正（如 7805 为 $+5\,\mathrm{V}$）
\\
79XX & GND & OUT & IN & 负（如 7905 为 $-5\,\mathrm{V}$）
\tabbot
\end{tabular}

\textbf{这是本册最容易出事的一张表}：78 与 79 的外形、封装、引脚数完全相同，
但 $1$ 号脚一个是 IN、一个是 GND。照 78 的接法去接 79，等于把电源与地反接，
上电瞬间就会烧。p38 的三个粉色提问框（\emph{What's the difference between 78 and 79?}、
\emph{What is inside the chip?}、\emph{What is the meaning of ××?}）正好对应三个知识点：

\begin{enumerate}
\item \textbf{78 与 79 的区别}：输出极性不同（正／负），因此管脚顺序不同（上表）。
\item \textbf{芯片里面是什么}：[教材补充] 基准电压源 + 取样分压 + 误差放大器 + 串联调整管 + 过流／过热保护。
从外部看，它就是一只「会自己调阻值的可变电阻」，串在输入与输出之间——这与稳压管的并联分流不同。
\item \textbf{×× 是什么意思}：额定输出电压的代号，$05\to5\,\mathrm{V}$、$12\to12\,\mathrm{V}$、$15\to15\,\mathrm{V}$、$24\to24\,\mathrm{V}$；
79XX 同规则但为负。
\end{enumerate}

\subhdr{8.6 压差与输入窗口（[教材补充]）}

p38 只讲了管脚，没有讲\textbf{这只器件需要什么输入}。而这一点是它能不能工作的前提，
也是它与稳压管电路最本质的差别：

\textbf{（1）输入必须比输出高出一个「压差」。} 三端稳压器内部是\textbf{串联}调整管，
它只能「降压」不能「升压」。内部调整管本身也要占去 $2\sim3\,\mathrm{V}$（称为压差，dropout），
所以输入有一个\textbf{下限}：
\[
U_{\mathrm{I,min}}\ge U_{\mathrm{O}}+U_{\mathrm{dropout}}+U_{\text{纹波}}/2
\quad\approx U_{\mathrm{O}}+(2\sim3)\,\mathrm{V}+\tfrac12\Delta U .
\]
对 7805，若滤波后纹波峰谷差 $3\,\mathrm{V}$，则输入至少要 $5+2+1.5=8.5\,\mathrm{V}$。
低于这个值，器件退出稳压，输出变成 $U_{\mathrm{O}}\approx U_{\mathrm{I}}-U_{\mathrm{dropout}}$，
\textbf{跟着输入一起往下掉}——这与 §6.3 的 (b)(e)、§8.3 的「输出等于输入」是同一种失效。

\textbf{（2）输入也不能太高，否则器件过热。} 输出电流全部穿过调整管，其耗散是
\[
P\approx(U_{\mathrm{I}}-U_{\mathrm{O}})\,I_{\mathrm{O}} .
\]
例如 7805 带 $0.5\,\mathrm{A}$ 负载，输入 $12\,\mathrm{V}$：$P=(12-5)\times0.5=3.5\,\mathrm{W}$，
必须加散热片；若输入取 $20\,\mathrm{V}$，同样负载下 $P=7.5\,\mathrm{W}$，多数塑封器件扛不住。
所以输入窗口是\textbf{两头都有限制}的：
\[
U_{\mathrm{O}}+U_{\mathrm{dropout}}\le U_{\mathrm{I}}\le U_{\mathrm{O}}+\frac{P_{\max}}{I_{\mathrm{O}}} .
\]
\textbf{（3）为什么工程上偏爱它。} 相比 §5 的稳压管电路，它的输出电流大得多（可达 $1\,\mathrm{A}$ 以上）、
不需要自己算限流电阻、内部还有过流与过热保护。代价是必须留压差、必须散热。

\aside{课件 p37--p38 只有分类图与管脚表，\textbf{没有提到压差、没有提功耗、没有提输入范围}。
本小节的式子与 $3.5\,\mathrm{W}$、$7.5\,\mathrm{W}$ 两个数都是本讲义补的（[教材补充]）。
补的理由：任务要求「三端稳压器的引脚与压差」，而压差是选型与故障判断的必要条件——
缺少它，学生只能背管脚，一旦电路不上电就无从下手。}

\subhdr{8.7 两条路线怎么分工：稳压管电路 vs 三端稳压器}

\par\vspace{2pt}
\noindent\begin{tabular}{@{}p{3.4cm}p{5.6cm}p{6.7cm}@{}}
\tabtop
\thead{比较项} & \thead{稳压管电路（3.6.1）} & \thead{三端稳压器（3.6.2）}
\\
\tabmid
连接方式 & 稳压管\textbf{并联}在负载两端，$R$ 串联分流 & 器件\textbf{串联}在输入与输出之间
\\
输出电流 & 小（多余电流都被稳压管吃掉） & 大（可达 $1\,\mathrm{A}$ 以上）
\\
效率 & 低（空载时电流全走稳压管） & 较高（但压差越大越热）
\\
外围元件 & 需要自己算 $R$（§5.2 的窗口） & 基本不用算，接三根线 + 输入输出电容
\\
失效方式 & 过流烧管、电流不足失稳 & 输入不足失稳、过热保护／烧毁
\\
\tabbot
\end{tabular}

一句话分工：\textbf{小电流、要便宜、当基准用稳压管；大电流、要省事、当电源用三端稳压器}。
两者在整机里是串联的——p14 那张成品板就是「整流滤波 → 三端稳压器」，而稳压管通常只出现在基准或保护支路里。

%=============================================================
\hdr{§9}{公式卡、四步清单与易错点}

\subhdr{9.1 公式卡（按做题顺序排，不按课件顺序）}

\par\vspace{2pt}{\small
\noindent\begin{tabular}{@{}p{4.3cm}p{6.3cm}p{4.9cm}@{}}
\tabtop
\thead{场合} & \thead{公式} & \thead{出处}
\\
\tabmid
桥式 + 电容的输出 & $U_{\mathrm{O}}=1.2U_2$,\quad $U_2=U_{\mathrm{O}}/1.2$ & p9、§2.1
\\
桥式 + 电容的耐压 & $U_{\mathrm{Rm}}=\sqrt2U_2$ & p9
\\
每管电流 & $I_{\mathrm{D}}=\tfrac12I_{\mathrm{O}}$（全波） & p9
\\
选管条件 & $I_{\mathrm{F}}\ge2I_{\mathrm{D}}$,\quad $U_{\mathrm{R}}\ge U_{\mathrm{Rm}}$ & p9
\\
选电容容量 & $C\ge\dfrac{1.5\sim2.5}{R_{\mathrm{L}}f}$，即 $R_{\mathrm{L}}C\ge(1.5\sim2.5)T$ & p10
\\
选电容耐压 & $U_{\mathrm{CN}}\ge\sqrt2U_2$ & p10
\\
两个极限 & 开路 $\sqrt2U_2$；断容 $0.9U_2$ & p11
\\
电感滤波纹波 & $\Delta U_{\mathrm{O}}/\Delta U_{\mathrm{in}}=R_{\mathrm{L}}/\sqrt{R_{\mathrm{L}}^2+(2\omega L)^2}$ & §2.5（补）
\\
$\pi$ 型 RC 的输出 & $U_{\mathrm{O}}\approx0.9U_2-I_{\mathrm{O}}R$ & §2.6（补）
\\
稳压管工作点 & $U_{\mathrm{O}}=U_{\mathrm{Z}}$,\ $I=(U_{\mathrm{I}}-U_{\mathrm{Z}})/R$,\ $I_{\mathrm{Z}}=I-U_{\mathrm{O}}/R_{\mathrm{L}}$ & p18、§5.1
\\
稳压判据 & $I_{\mathrm{Zmin}}\le I_{\mathrm{Z}}\le I_{\mathrm{Zmax}}$ & p18
\\
调节率 & $\Delta U_{\mathrm{O}}/\Delta U_{\mathrm{I}}=r_{\mathrm{Z}}/(R+r_{\mathrm{Z}})$ & §3.3（补）
\\
输出电阻 & $\Delta U_{\mathrm{O}}/\Delta I_{\mathrm{O}}=-(R\parallel r_{\mathrm{Z}})$ & §3.3（补）
\\
$R$ 的窗口 & $\dfrac{U_{\mathrm{Imax}}-U_{\mathrm{Z}}}{I_{\mathrm{Zmax}}+I_{\mathrm{Omin}}}\le R\le\dfrac{U_{\mathrm{Imin}}-U_{\mathrm{Z}}}{I_{\mathrm{Zmin}}+I_{\mathrm{Omax}}}$ & §5.2（补）
\\
三端稳压器 & $U_{\mathrm{I}}\ge U_{\mathrm{O}}+(2\sim3)\,\mathrm{V}$,\quad $P\approx(U_{\mathrm{I}}-U_{\mathrm{O}})I_{\mathrm{O}}$ & §8.6（补）
\tabbot
\end{tabular}}

\subhdr{9.2 四步清单（拿到任何稳压题都照这个顺序）}

\begin{enumerate}
\item \textbf{定结构}：输入极性正还是反？管子是串联还是并联？输出取在哪两点之间？（§6.1 的表）
\item \textbf{判区}：反接且两管反向时，先比 $U_{\mathrm{I}}$ 与 $U_{\mathrm{Z1}}+U_{\mathrm{Z2}}$；
反接且管子正向时，压降直接取 $2U_{\mathrm{F}}$；输入反号时，别忘了 $U_{\mathrm{O}}$ 也跟着反号。
\item \textbf{代公式}：只在\textbf{击穿区成立}时才用 $U_{\mathrm{O}}=U_{\mathrm{Z}}$，其余情况一律走 $U_{\mathrm{O}}\approx U_{\mathrm{I}}$ 或正向压降。
\item \textbf{验算}：$I_{\mathrm{Z}}$ 是否在 $I_{\mathrm{Zmin}}\sim I_{\mathrm{Zmax}}$ 之间、$R$ 与管子的功率够不够。
\end{enumerate}

\subhdr{9.3 易错点（每条都能指回某页某步）}

\begin{enumerate}
\item \textbf{$34.5$ 与 $35.4$}（p9）：同一页两个数自相矛盾，正确的是 $35.4\,\mathrm{V}$。
\item \textbf{$I_{\mathrm{F}}$ 乘 $2$、$U_{\mathrm{R}}$ 不乘}（p9）：一个看有效值、一个看峰值。
\item \textbf{半波带电容的 $2\sqrt2U_2$}（p1）：只对半波成立，桥式带电容仍是 $\sqrt2U_2$。
\item \textbf{「$U_{\mathrm{O}}=1.2U_2$」不是常数}（p10、§2.4）：它随 $C$ 与 $R_{\mathrm{L}}$ 变，负载开路时到 $\sqrt2U_2$。
\item \textbf{把 $\pi$ 型 RC 的串联元件当电感}（p13）：图里画的是线圈，按标题应为 $R$。
\item \textbf{反接就以为能稳压}（p23、p24）：反接 = 正向导通 = 只有 $1\,\mathrm{V}$ 压降。
\item \textbf{输入不足却仍写 $U_{\mathrm{O}}=U_{\mathrm{Z}}$}（p22、p25）：此时 $U_{\mathrm{O}}\approx U_{\mathrm{I}}$，方程无解（§7.1）。
\item \textbf{只算 $I_{\mathrm{Z}}$ 不判限}（p27）：本题答案落在「超限」这一档（§7.2）。
\item \textbf{两条支路相减当稳压电源}（p26）：误差相加（$\pm14\%$）、电流各自受限（§7.3）。
\item \textbf{78 与 79 照抄接法}（p38）：$1$ 脚一个是 IN、一个是 GND，接反即烧（§8.5）。
\item \textbf{忘了压差}（p38）：三端稳压器只能降压，输入低于 $U_{\mathrm{O}}+2\,\mathrm{V}$ 就失稳（§8.6）。
\end{enumerate}

\subhdr{9.4 量级自检：一秒钟判断算没算错}

\begin{itemize}
\item 变压器二次电压 $U_2$ 通常是几伏到几十伏；桥式带电容的 $U_{\mathrm{O}}$ 只能是 $1.2U_2$ 左右，
\textbf{不可能超过 $\sqrt2U_2$}（那是负载开路的上限）。
\item 稳压管的 $I_{\mathrm{Z}}$ 通常是几毫安到几十毫安；若算出几百毫安，一定是 $R$ 取小了或单位错了。
\item 限流电阻 $R$ 通常是几十到几百欧；若算出几欧或几十千欧，回去检查 $I_{\mathrm{Zmax}}$ 与 $R_{\mathrm{L}}$。
\item 三端稳压器的压差是 $2\sim3\,\mathrm{V}$ 量级；若设计的输入只比输出高 $0.5\,\mathrm{V}$，方案不成立。
\end{itemize}

%=============================================================
\hdr{§10}{缺口清单：课件哪句话跳了步，本讲义补在哪}

共 14 条，按课件页码排序。每条五个字段：类型｜页码｜课件原文短语｜缺了什么｜补法（写到哪一节）。

\gapitem{机制缺口}{p20\quad\emph{``$U_{\mathrm{O}}\uparrow\rightarrow I_{\mathrm{Z}}\uparrow\uparrow\rightarrow I\uparrow
\rightarrow IR\uparrow\cdots U_{\mathrm{O}}\downarrow$''}\newline
缺：五步里每一跳的出处（全页无一个基尔霍夫方程），以及「$\uparrow\uparrow$」那两个箭头到底多快。
\newline 补：§4.2 把五跳分别对上 $U_{\mathrm{O}}=U_{\mathrm{Z}}+r_{\mathrm{Z}}I_{\mathrm{Z}}$、
$I=I_{\mathrm{Z}}+I_{\mathrm{O}}$、$U_{\mathrm{O}}=U_{\mathrm{I}}-IR$；§3.3 给出倍数 $r_{\mathrm{Z}}/(R+r_{\mathrm{Z}})$。}

\gapitem{适用区间}{p18\quad\emph{``$I_{\mathrm{Zmin}}$ \ldots There is no regulating effect''}／
\emph{``$I_{\mathrm{Zmax}}$ \ldots It will be damaged forever''}\newline
缺：两句只定义了边界，没给「给定 $U_{\mathrm{I}}$、$R$、$R_{\mathrm{L}}$，怎么验算 $I_{\mathrm{Z}}$ 落不落在里面」。
\newline 补：§5.1 三步法（先假定 $U_{\mathrm{O}}=U_{\mathrm{Z}}$，再算 $I_{\mathrm{Z}}$，再分三档判定）。}

\gapitem{适用区间}{p18\quad\emph{``$r_{\mathrm{Z}}=\Delta U/\Delta I$ \ldots The smaller the better''}\newline
缺：只给定义与一句形容词，没给「好在哪里」的定量结果，也没说明该式只在击穿段成立。
\newline 补：§3.3 两条用法——输入扰动的调节率 $r_{\mathrm{Z}}/(R+r_{\mathrm{Z}})$、输出电阻 $-(R\parallel r_{\mathrm{Z}})$。}

\gapitem{代数跳步}{p9\quad\emph{``$U_{\mathrm{R}}\ge U_{\mathrm{Rm}}=34.5\,\mathrm{V}$''}\newline
缺：同一页上一行刚给 $U_{\mathrm{Rm}}=\sqrt2\times25\,\mathrm{V}=35.4\,\mathrm{V}$，两数自相矛盾。
\newline 补：§2.2 点名这是数字转置，全文按 $35.4\,\mathrm{V}$ 走。}

\gapitem{机制缺口}{p9\quad\emph{``$I_{\mathrm{F}}\ge2I_{\mathrm{D}}=2\times0.15\,\mathrm{A}$''}\newline
缺：为什么乘 $2$、而 $U_{\mathrm{R}}\ge U_{\mathrm{Rm}}$ 不乘，课件不解释。
\newline 补：§2.2 指出一个是有效值（窄脉冲波形因数）、一个是峰值。}

\gapitem{机制缺口}{p10\quad\emph{``$C\ge\dfrac{1.5\sim2.5}{R_{\mathrm{L}}f}$''}\newline
缺：经验式的来历；为什么是 $1.5\sim2.5$、为什么除以 $f$。
\newline 补：§2.3 改写为等价形式 $R_{\mathrm{L}}C\ge(1.5\sim2.5)T$，并给出「撑到下一个峰值」的物理含义。}

\gapitem{拼接缺口}{p11\quad\emph{``when the load is open circuit''}／\emph{``when the capacitor branch is open''}\newline
缺：两个极限算完就停，没有一句话把它们接到 3.6 的稳压级上（$35.4$ 与 $22.5\,\mathrm{V}$ 差 $57\%$ 这个事实未被利用）。
\newline 补：§1.2 与 §2.4 两次落这个缝；§1.1 给出四段链条的分工。}

\gapitem{机制缺口}{p12\quad\emph{``Normally used in the situation when the output voltage is relatively lower
while the current through the load is relatively larger.''}\newline
缺：为什么「低压大电流」场合该用电感，课件只有结论。
\newline 补：§2.5 说明电感不抬高直流（$U_{\mathrm{O}}\approx0.9U_2$）只衰减纹波，并给出 $R_{\mathrm{L}}/\sqrt{R_{\mathrm{L}}^2+(2\omega L)^2}$ 的数值表。}

\gapitem{形式改写}{p13\quad\emph{``（2）$\pi$-type RC Filter''}，而图内串联元件画成线圈并标 \emph{L}，上方另有红字 \emph{R}\newline
缺：标题说 RC，器件画成 L，两者矛盾，未说明以哪个为准。
\newline 补：§2.6 按标题取 $R$，并给出 RC 版的代价 $U_{\mathrm{O}}\approx0.9U_2-I_{\mathrm{O}}R$ 与适用电流。}

\gapitem{拼接缺口}{p16、p21\quad\emph{``designed to operate in the reverse breakdown region''}／
\emph{``the zener diode need to be connected in reverse''}\newline
缺：「击穿为什么可逆、为什么不损坏」只有断言；从「器件参数」到「电路判据」之间也没有桥。
\newline 补：§3.1 把击穿的可逆性落到「限流」这一个条件上；§5 补参数到判据的三步法。}

\gapitem{答案形态}{p22、p25\quad（无原文，只有电路与 $U_{\mathrm{I}}=10\,\mathrm{V}$／$8\,\mathrm{V}$）\newline
缺：输出为什么只能写约等号、为什么算不出精确值，课件不交代。
\newline 补：§7.1 说明此时伏安特性在击穿前近似开路，精确值需要漏电阻而该参数课件没给，故答案停在 $U_{\mathrm{O}}\approx U_{\mathrm{I}}$。}

\gapitem{机制缺口}{p27\quad\emph{``（g）$U_{\mathrm{Z}}=10\,\mathrm{V}$\ $I_{\mathrm{Zmax}}=23\,\mathrm{mA}$. Find $I_{\mathrm{Z}}$''}\newline
缺：只问电流，没给「先假定再验算」的流程，也没说超限后该怎么办。
\newline 补：§6.7 全程三步法，$I_{\mathrm{Z}}=26.7\,\mathrm{mA}$ 判超限，并反解出 $R\ge337\,\Omega$ 的改法；§7.2 说明这类题的答案形态。}

\gapitem{拼接缺口}{p28--p36\quad\emph{``4. Example 2''}，万用表读数 \emph{000}\newline
缺：九帧同一张图，无数据、无解题过程、无故障分析；页眉章节号在 3.6 与 8.6 之间跳动。
\newline 补：§8.1 去重说明并点名编号笔误；§8.2 给出整机五步关系链；§8.3 给出读数 $000$ 的四行诊断表。}

\gapitem{机制缺口}{p37--p38\quad\emph{``Classification''}、\emph{``Shape and pins''}、
\emph{``What's the difference between 78 and 79?''}、\emph{``What is the meaning of ××?''}\newline
缺：三个提问都没有答案；管脚之外没有压差、没有输入范围、没有功耗，于是「照管脚接上却不上电」时无从判断。
\newline 补：§8.5 逐条答三问（含 $XX$ 的含义）；§8.6 补压差与输入窗口 $U_{\mathrm{I}}\ge U_{\mathrm{O}}+(2\sim3)\,\mathrm{V}$ 与 $P\approx(U_{\mathrm{I}}-U_{\mathrm{O}})I_{\mathrm{O}}$。}

\subhdr{10.4 本册补的四处（课件完全没有的内容）}

\par\vspace{2pt}
\noindent\begin{tabular}{@{}p{3.2cm}p{1.7cm}p{11.0cm}@{}}
\tabtop
\thead{补的内容} & \thead{对应页} & \thead{为什么必须补}
\\
\tabmid
$r_{\mathrm{Z}}$ 的两个定量用法 & p18 & 课件只说「越小越好」；不给比值，就无法回答「这个电源稳不稳」、也算不了 (f) 的精度
\\
$R$ 的取值窗口 & p18、p27 & 课件只给判据不给设计式；而考试与工程问的都是「$R$ 取多少」
\\
三端稳压器的压差与功耗 & p38 & 只背管脚接不出能工作的电路；压差决定最低输入，功耗决定散热
\\
Example 2 的故障诊断表 & p28--p36 & 一张只有读数的图，不教怎么读那个数，就没有信息量
\tabbot
\end{tabular}

%=============================================================
\hdr{§11}{三句话总结}

\textbf{第一句}：滤波与稳压解决的是两件事——滤波把脉动削平但留下了「随负载漂的输出」（$30\to35.4\to22.5\,\mathrm{V}$），
稳压把输出钉在 $U_{\mathrm{Z}}$ 上，但代价是\textbf{必须让稳压管吃掉多余的电压}。

\textbf{第二句}：稳压管电路的一切判断只用三个方程加一条不等式——
$U_{\mathrm{O}}=U_{\mathrm{Z}}+r_{\mathrm{Z}}I_{\mathrm{Z}}$、$I=I_{\mathrm{Z}}+I_{\mathrm{O}}$、$U_{\mathrm{O}}=U_{\mathrm{I}}-IR$，
以及 $I_{\mathrm{Zmin}}\le I_{\mathrm{Z}}\le I_{\mathrm{Zmax}}$；Example 1 的七个小问，全是在这条不等式两侧取边界。

\textbf{第三句}：三端稳压器把「并联分流」换成了「串联降压」，于是不再需要自己算限流电阻，
代价是多出一条\textbf{压差}约束（$U_{\mathrm{I}}\ge U_{\mathrm{O}}+2\sim3\,\mathrm{V}$）与一个散热问题——
管脚（78：IN--GND--OUT；79：GND--OUT--IN）接错或输入不够，是这一节最常见的两种故障。

%=============================================================
\hdr{§12}{与课件的页码对照表}

\par\vspace{2pt}{\footnotesize
\noindent\begin{tabular}{@{}p{1.7cm}p{8.6cm}p{5.6cm}@{}}
\tabtop
\thead{课件页} & \thead{内容（本册的事实来源）} & \thead{本讲义}
\\
\tabmid
p1 & 两道选择题 + 四列参数总表 & §1.3、§2.1
\\
p2、p3 & Review：半波 $U_{\mathrm{Rm}}=\sqrt2U_2$、中心抽头 $2\sqrt2U_2$ & §1.3
\\
p4、p5 & Review：桥式 $U_{\mathrm{Rm}}=\sqrt2U_2$（含导通路径） & §1.3、§1.2
\\
p6 & 直流稳压电源实物（适配器） & §0.1
\\
p7 & 两套 LED 灯串电路 + 实物照片 & §1.4
\\
p8 & 3.5.2 电容滤波：E1 题面 + 总表 & §2.1
\\
p9 & E1（1）选二极管（含 $34.5$ 笔误） & §2.2
\\
p10 & E1（2）选电容（容量与耐压） & §2.3
\\
p11 & E1（3）（4）负载开路／电容支路断开 & §2.4、§1.2
\\
p12 & 3.5.3 电感滤波 & §2.5
\\
p13 & 3.5.4 复式滤波（$\pi$ 型 LC／RC） & §2.6
\\
p14 & 3.6 路标：稳压管、三端稳压器实物 & §2.7
\\
p15 & 稳压管与普通二极管的伏安特性对比 & §3.1
\\
p16 & 稳压管定义（反接、击穿区工作） & §3.1
\\
p17 & 特性两条 + 五类应用 & §3.1
\\
p18 & 2. Parameters：$U_{\mathrm{Z}}$、$I_{\mathrm{Zmin}}$、$I_{\mathrm{Zmax}}$、$r_{\mathrm{Z}}$ & §3.2、§3.3
\\
p19 & 3. 稳压管稳压电路（整机图） & §4.1
\tabbot
\end{tabular}}

\par\vspace{2pt}{\footnotesize
\noindent\begin{tabular}{@{}p{1.7cm}p{8.6cm}p{5.6cm}@{}}
\tabtop
\thead{课件页（续）} & \thead{内容（本册的事实来源）} & \thead{本讲义}
\\
\tabmid
p20 & 稳压过程箭头链（$I$、$I_{\mathrm{Z}}$、$I_{\mathrm{O}}$ 方向） & §4.2、§4.3
\\
p21 & Example 1 (a)：$U_{\mathrm{I}}=18\,\mathrm{V}$ + 两个 Notice & §6.1、§6.2
\\
p22 & (b)：$U_{\mathrm{I}}=10\,\mathrm{V}$ & §6.3、§7.1
\\
p23 & (c)：$U_{\mathrm{I}}=4\,\mathrm{V}$，输入反接 & §6.4
\\
p24 & (d)：$U_{\mathrm{I}}=4\,\mathrm{V}$，输入反接（与 (c) 标注相同，[待确认]） & §6.4
\\
p25 & (e)：$U_{\mathrm{I}}=8\,\mathrm{V}$ & §6.5
\\
p26 & (f)：$U_{\mathrm{I}}=13\,\mathrm{V}$，两支路相减 & §6.6、§7.3
\\
p27 & (g)：$U_{\mathrm{I}}=20\,\mathrm{V}$、$R=300\,\Omega$、$R_{\mathrm{L}}=1.5\,\mathrm{k\Omega}$ & §6.7、§7.2
\\
p28--p36 & Example 2：整机电路 + 万用表读数 $000$（九帧） & §8.1--§8.3
\\
p37 & 3.6.2 集成稳压器分类图 & §8.4
\\
p38 & 三端稳压器外形与管脚（7805／7905）+ 三个提问 & §8.5、§8.6
\tabbot
\end{tabular}}

\flowerrule

{\small\color{inkgray}
算术用 Python 复算：$U_2=25\,\mathrm{V}$、$U_{\mathrm{Rm}}=35.4\,\mathrm{V}$、$14=4.5+9.5$、
$I_{\mathrm{Z}}=26.67\,\mathrm{mA}$、$R\ge337\,\Omega$、$P=3.5\,\mathrm{W}$；
课件算错的 $34.5\,\mathrm{V}$ 已在 §2.2 点名并按 $35.4\,\mathrm{V}$ 走。}

\end{document}
